\documentclass[10.5pt]{article}
\usepackage[letterpaper,margin=0.72in]{geometry}
\usepackage{booktabs,longtable,array,tabularx}
\usepackage{xcolor,hyperref,xurl,fancyhdr,lastpage,microtype}
\usepackage[T1]{fontenc}\usepackage{lmodern}
\definecolor{navy}{RGB}{25,54,93}\definecolor{light}{RGB}{242,245,249}
\hypersetup{colorlinks=true,linkcolor=navy,urlcolor=navy}
\setlength{\parindent}{0pt}\setlength{\parskip}{0.38em}
\pagestyle{fancy}\fancyhf{}\lhead{\small FMCW / TWTT Research Corpus}\rhead{\small Annotated Source Atlas}\cfoot{\small Page \thepage\ of \pageref{LastPage}}
\begin{document}
\begin{titlepage}\centering\vspace*{1.2in}{\Huge\bfseries\color{navy} FMCW/TWTT Annotated Source Atlas\par}\vspace{0.3in}{\Large 84-source map for the Saeed AWR2944 timing project\par}\vspace{0.5in}{\large Summer 2026\par}\vfill\begin{minipage}{0.84\textwidth}\small This atlas is the source-by-source companion to the literature review, mathematical handbook, and MATLAB simulation blueprint. It records why each source is in the corpus, what to extract from it, whether a physical research copy is present, and where it lives in the archive. It is deliberately exhaustive so the literature repository remains navigable after the first simulation sprint.\end{minipage}\end{titlepage}
\tableofcontents\clearpage
\section{How to use this atlas}
The atlas is not a substitute for reading the priority papers. It is a navigation and extraction map. For every paper that becomes active in the project, record (1) waveform/architecture, (2) unknown parameters, (3) estimator, (4) required synchronization/coherence assumptions, (5) demonstrated precision/accuracy, and (6) one limitation that matters for two AWR2944 boards.

The priority labels are working labels: P0 is foundational/immediate, P1 is strong second-wave material, P2 is useful supporting literature, and ``reference'' denotes documentation or context. ``Citation only'' means the source remains in the research map but its PDF bytes are not present in the current archive.
\clearpage
\section{P0 - 77 GHz cooperative FMCW}

\subsection*{1. A 77-GHz Cooperative Radar System Based on Multi-Channel FMCW Stations for Local Positioning Applications (2013)}
\begin{tabularx}{\textwidth}{@{}p{0.18\textwidth}X@{}}
\textbf{Priority/status} & P0 / \nolinkurl{citation_only} \\
\textbf{Authors} & R. Feger et al. \\
\textbf{DOI / identifier} & \nolinkurl{10.1109/TMTT.2012.2227781} \\
\textbf{Archive path} & not present \\
\end{tabularx}
\textbf{Why it is here.} Important bridge to 77-GHz hardware.

\textbf{What to extract for this project.} Use to understand which observables survive separate oscillators, what degree of coherence is required, and how cross-path data should be retained/combined across nodes.

\textbf{Notebook prompt.} Write down the paper's assumptions in equation form. Mark each assumption as \emph{true in V0}, \emph{true only after hardware synchronization}, \emph{requires calibration}, or \emph{not true for two independent AWR boards}. This classification is more valuable than a generic summary.

\vspace{0.35em}\hrule\vspace{0.65em}

\clearpage
\section{P0 - FMCW fundamentals}

\subsection*{2. Linear FMCW Radar Techniques (1992)}
\begin{tabularx}{\textwidth}{@{}p{0.18\textwidth}X@{}}
\textbf{Priority/status} & P0 / \nolinkurl{citation_only} \\
\textbf{Authors} & A. G. Stove \\
\textbf{DOI / identifier} & \nolinkurl{10.1049/ip-f-2.1992.0048} \\
\textbf{Archive path} & not present \\
\end{tabularx}
\textbf{Why it is here.} Classic reference.

\textbf{What to extract for this project.} Reference for FMCW/radar fundamentals or adjacent methods. Use it to sanity-check notation, signal relationships, and assumptions when its topic becomes active in the implementation.

\textbf{Notebook prompt.} Write down the paper's assumptions in equation form. Mark each assumption as \emph{true in V0}, \emph{true only after hardware synchronization}, \emph{requires calibration}, or \emph{not true for two independent AWR boards}. This classification is more valuable than a generic summary.

\vspace{0.35em}\hrule\vspace{0.65em}

\clearpage
\section{P0 - MATLAB / simulation}

\subsection*{3. A Versatile FMCW Radar System Simulator for Millimeter-Wave Applications (2008)}
\begin{tabularx}{\textwidth}{@{}p{0.18\textwidth}X@{}}
\textbf{Priority/status} & P0 / \nolinkurl{included_restricted_user_supplied} \\
\textbf{Authors} & S. Scheiblhofer; M. Treml; S. Schuster; R. Feger; A. Stelzer \\
\textbf{DOI / identifier} & \nolinkurl{10.1109/EUMC.2008.4751778} \\
\textbf{Archive path} & \path{09_RESTRICTED_OR_CITATION_ONLY/user_supplied_restricted/2008_Scheiblhofer_A_Versatile_FMCW_Radar_System_Simulator.pdf} \\
\end{tabularx}
\textbf{Why it is here.} Modular MATLAB simulator from sweep synthesis through baseband.

\textbf{What to extract for this project.} Use to define simulator architecture, signal-domain fidelity, validation strategy, or efficient baseband implementation. Extract module boundaries, required state variables, and validation metrics.

\textbf{Notebook prompt.} Write down the paper's assumptions in equation form. Mark each assumption as \emph{true in V0}, \emph{true only after hardware synchronization}, \emph{requires calibration}, or \emph{not true for two independent AWR boards}. This classification is more valuable than a generic summary.

\vspace{0.35em}\hrule\vspace{0.65em}

\clearpage
\section{P0 - TWTT / cooperative FMCW}

\subsection*{4. Method for High Precision Radar Distance Measurement and Synchronization of Wireless Units (2007)}
\begin{tabularx}{\textwidth}{@{}p{0.18\textwidth}X@{}}
\textbf{Priority/status} & P0 / \nolinkurl{included_restricted_user_supplied} \\
\textbf{Authors} & S. Roehr; M. Vossiek; P. Gulden \\
\textbf{DOI / identifier} & \nolinkurl{10.1109/MWSYM.2007.380436} \\
\textbf{Archive path} & \path{09_RESTRICTED_OR_CITATION_ONLY/user_supplied_restricted/2007_Roehr_High_Precision_Radar_Distance_Synchronization.pdf} \\
\end{tabularx}
\textbf{Why it is here.} Core predecessor; <100 ps time synchronization reported.

\textbf{What to extract for this project.} Use to formalize clock epoch offset, clock-rate/frequency offset, reciprocity, two-way cancellation, presynchronization, and calibration. Translate each distinct clock/path error into an explicit simulator parameter.

\textbf{Notebook prompt.} Write down the paper's assumptions in equation form. Mark each assumption as \emph{true in V0}, \emph{true only after hardware synchronization}, \emph{requires calibration}, or \emph{not true for two independent AWR boards}. This classification is more valuable than a generic summary.

\vspace{0.35em}\hrule\vspace{0.65em}

\subsection*{5. Coherent Full-Duplex Double-Sided Two-Way Ranging and Velocity Measurement Between Separate Incoherent Radio Units (2019)}
\begin{tabularx}{\textwidth}{@{}p{0.18\textwidth}X@{}}
\textbf{Priority/status} & P0 / \nolinkurl{included_restricted_user_supplied} \\
\textbf{Authors} & M. Gottinger; F. Kirsch; P. Gulden; M. Vossiek \\
\textbf{DOI / identifier} & \nolinkurl{10.1109/TMTT.2019.2902553} \\
\textbf{Archive path} & \path{09_RESTRICTED_OR_CITATION_ONLY/user_supplied_restricted/2019_Gottinger_CFDDS_TWR.pdf} \\
\end{tabularx}
\textbf{Why it is here.} Closest academic analogue: reciprocal FMCW with separate clocks.

\textbf{What to extract for this project.} Use to formalize clock epoch offset, clock-rate/frequency offset, reciprocity, two-way cancellation, presynchronization, and calibration. Translate each distinct clock/path error into an explicit simulator parameter.

\textbf{Notebook prompt.} Write down the paper's assumptions in equation form. Mark each assumption as \emph{true in V0}, \emph{true only after hardware synchronization}, \emph{requires calibration}, or \emph{not true for two independent AWR boards}. This classification is more valuable than a generic summary.

\vspace{0.35em}\hrule\vspace{0.65em}

\clearpage
\section{P0 - clock synchronization}

\subsection*{6. Method for High Precision Clock Synchronization in Wireless Systems with Application to Radio Navigation (2007)}
\begin{tabularx}{\textwidth}{@{}p{0.18\textwidth}X@{}}
\textbf{Priority/status} & P0 / \nolinkurl{citation_only} \\
\textbf{Authors} & S. Roehr; P. Gulden; M. Vossiek \\
\textbf{DOI / identifier} & \nolinkurl{10.1109/RWS.2007.351890} \\
\textbf{Archive path} & not present \\
\end{tabularx}
\textbf{Why it is here.} Foundational clock synchronization; full copy not physically present in this runtime.

\textbf{What to extract for this project.} Use to formalize clock epoch offset, clock-rate/frequency offset, reciprocity, two-way cancellation, presynchronization, and calibration. Translate each distinct clock/path error into an explicit simulator parameter.

\textbf{Notebook prompt.} Write down the paper's assumptions in equation form. Mark each assumption as \emph{true in V0}, \emph{true only after hardware synchronization}, \emph{requires calibration}, or \emph{not true for two independent AWR boards}. This classification is more valuable than a generic summary.

\vspace{0.35em}\hrule\vspace{0.65em}

\clearpage
\section{P0 - cooperative FMCW}

\subsection*{7. Precise Distance Measurement with Cooperative FMCW Radar Units (2008)}
\begin{tabularx}{\textwidth}{@{}p{0.18\textwidth}X@{}}
\textbf{Priority/status} & P0 / \nolinkurl{citation_only} \\
\textbf{Authors} & A. Stelzer; M. Jahn; S. Scheiblhofer \\
\textbf{DOI / identifier} & \nolinkurl{10.1109/RWS.2008.4463606} \\
\textbf{Archive path} & not present \\
\end{tabularx}
\textbf{Why it is here.} Core cooperative FMCW ranging paper.

\textbf{What to extract for this project.} Use to understand which observables survive separate oscillators, what degree of coherence is required, and how cross-path data should be retained/combined across nodes.

\textbf{Notebook prompt.} Write down the paper's assumptions in equation form. Mark each assumption as \emph{true in V0}, \emph{true only after hardware synchronization}, \emph{requires calibration}, or \emph{not true for two independent AWR boards}. This classification is more valuable than a generic summary.

\vspace{0.35em}\hrule\vspace{0.65em}

\subsection*{8. Precise Distance and Velocity Measurement for Real Time Locating in Multipath Environments Using a Frequency-Modulated Continuous-Wave Secondary Radar Approach (2008)}
\begin{tabularx}{\textwidth}{@{}p{0.18\textwidth}X@{}}
\textbf{Priority/status} & P0 / \nolinkurl{citation_only} \\
\textbf{Authors} & S. Roehr; P. Gulden; M. Vossiek \\
\textbf{DOI / identifier} & \nolinkurl{10.1109/TMTT.2008.2003137} \\
\textbf{Archive path} & not present \\
\end{tabularx}
\textbf{Why it is here.} Detailed range/velocity/multipath model.

\textbf{What to extract for this project.} Use to understand which observables survive separate oscillators, what degree of coherence is required, and how cross-path data should be retained/combined across nodes.

\textbf{Notebook prompt.} Write down the paper's assumptions in equation form. Mark each assumption as \emph{true in V0}, \emph{true only after hardware synchronization}, \emph{requires calibration}, or \emph{not true for two independent AWR boards}. This classification is more valuable than a generic summary.

\vspace{0.35em}\hrule\vspace{0.65em}

\clearpage
\section{P0 - frequency/phase estimation}

\subsection*{9. Accuracy Limits of a K-Band FMCW Radar with Phase Evaluation (2012)}
\begin{tabularx}{\textwidth}{@{}p{0.18\textwidth}X@{}}
\textbf{Priority/status} & P0 / \nolinkurl{citation_only} \\
\textbf{Authors} & S. Scherr; S. Ayhan; M. Pauli; T. Zwick \\
\textbf{DOI / identifier} & not recorded \\
\textbf{Archive path} & not present \\
\end{tabularx}
\textbf{Why it is here.} CRLB and phase-evaluation precision precedent.

\textbf{What to extract for this project.} Use for the precision/error-budget layer. Implement the impairment before the estimator whenever possible, then compare bias/variance/RMSE with the paper's analytical or measured behavior.

\textbf{Notebook prompt.} Write down the paper's assumptions in equation form. Mark each assumption as \emph{true in V0}, \emph{true only after hardware synchronization}, \emph{requires calibration}, or \emph{not true for two independent AWR boards}. This classification is more valuable than a generic summary.

\vspace{0.35em}\hrule\vspace{0.65em}

\subsection*{10. An Efficient Frequency and Phase Estimation Algorithm With CRB Performance for FMCW Radar Applications (2015)}
\begin{tabularx}{\textwidth}{@{}p{0.18\textwidth}X@{}}
\textbf{Priority/status} & P0 / \nolinkurl{included_restricted_user_supplied} \\
\textbf{Authors} & S. Scherr et al. \\
\textbf{DOI / identifier} & \nolinkurl{10.1109/TIM.2014.2381354} \\
\textbf{Archive path} & \path{09_RESTRICTED_OR_CITATION_ONLY/user_supplied_restricted/2015_Scherr_Efficient_Frequency_and_Phase_Estimation_CRB.pdf} \\
\end{tabularx}
\textbf{Why it is here.} CZT plus phase evaluation; CRB context.

\textbf{What to extract for this project.} Use for the precision/error-budget layer. Implement the impairment before the estimator whenever possible, then compare bias/variance/RMSE with the paper's analytical or measured behavior.

\textbf{Notebook prompt.} Write down the paper's assumptions in equation form. Mark each assumption as \emph{true in V0}, \emph{true only after hardware synchronization}, \emph{requires calibration}, or \emph{not true for two independent AWR boards}. This classification is more valuable than a generic summary.

\vspace{0.35em}\hrule\vspace{0.65em}

\clearpage
\section{P0 - multistatic synchronization}

\subsection*{11. FMCW Radar Transceiver Synchronization in a Multistatic Microwave Tomography System (2023)}
\begin{tabularx}{\textwidth}{@{}p{0.18\textwidth}X@{}}
\textbf{Priority/status} & P0 / \nolinkurl{citation_only} \\
\textbf{Authors} & B. Lettner \\
\textbf{DOI / identifier} & not recorded \\
\textbf{Archive path} & not present \\
\end{tabularx}
\textbf{Why it is here.} Thesis-length synchronization treatment.

\textbf{What to extract for this project.} Use to formalize clock epoch offset, clock-rate/frequency offset, reciprocity, two-way cancellation, presynchronization, and calibration. Translate each distinct clock/path error into an explicit simulator parameter.

\textbf{Notebook prompt.} Write down the paper's assumptions in equation form. Mark each assumption as \emph{true in V0}, \emph{true only after hardware synchronization}, \emph{requires calibration}, or \emph{not true for two independent AWR boards}. This classification is more valuable than a generic summary.

\vspace{0.35em}\hrule\vspace{0.65em}

\clearpage
\section{P0 - nonidealities}

\subsection*{12. Impact of Frequency Ramp Nonlinearity, Phase Noise, and SNR on FMCW Radar Accuracy (2016)}
\begin{tabularx}{\textwidth}{@{}p{0.18\textwidth}X@{}}
\textbf{Priority/status} & P0 / \nolinkurl{included_restricted_user_supplied} \\
\textbf{Authors} & S. Ayhan et al. \\
\textbf{DOI / identifier} & \nolinkurl{10.1109/TMTT.2016.2599165} \\
\textbf{Archive path} & \path{09_RESTRICTED_OR_CITATION_ONLY/user_supplied_restricted/2016_Ayhan_Ramp_Nonlinearity_Phase_Noise_SNR.pdf} \\
\end{tabularx}
\textbf{Why it is here.} Analytical and measured accuracy effects.

\textbf{What to extract for this project.} Use for the precision/error-budget layer. Implement the impairment before the estimator whenever possible, then compare bias/variance/RMSE with the paper's analytical or measured behavior.

\textbf{Notebook prompt.} Write down the paper's assumptions in equation form. Mark each assumption as \emph{true in V0}, \emph{true only after hardware synchronization}, \emph{requires calibration}, or \emph{not true for two independent AWR boards}. This classification is more valuable than a generic summary.

\vspace{0.35em}\hrule\vspace{0.65em}

\clearpage
\section{P0 - phase noise}

\subsection*{13. The Effect of Phase Noise on Ranging Uncertainty in FMCW Secondary Radar-Based Local Positioning Systems (2012)}
\begin{tabularx}{\textwidth}{@{}p{0.18\textwidth}X@{}}
\textbf{Priority/status} & P0 / \nolinkurl{citation_only} \\
\textbf{Authors} & R. Ebelt; D. Shmakov; M. Vossiek \\
\textbf{DOI / identifier} & not recorded \\
\textbf{Archive path} & not present \\
\end{tabularx}
\textbf{Why it is here.} Directly relevant to secondary/cooperative FMCW PN.

\textbf{What to extract for this project.} Use for the precision/error-budget layer. Implement the impairment before the estimator whenever possible, then compare bias/variance/RMSE with the paper's analytical or measured behavior.

\textbf{Notebook prompt.} Write down the paper's assumptions in equation form. Mark each assumption as \emph{true in V0}, \emph{true only after hardware synchronization}, \emph{requires calibration}, or \emph{not true for two independent AWR boards}. This classification is more valuable than a generic summary.

\vspace{0.35em}\hrule\vspace{0.65em}

\clearpage
\section{P0 - phase noise / PLL}

\subsection*{14. Analysis of Ranging Precision in an FMCW Radar Measurement Using a Phase-Locked Loop (2018)}
\begin{tabularx}{\textwidth}{@{}p{0.18\textwidth}X@{}}
\textbf{Priority/status} & P0 / \nolinkurl{included_restricted_user_supplied} \\
\textbf{Authors} & F. Herzel; D. Kissinger; H. J. Ng \\
\textbf{DOI / identifier} & \nolinkurl{10.1109/TCSI.2017.2733041} \\
\textbf{Archive path} & \path{09_RESTRICTED_OR_CITATION_ONLY/user_supplied_restricted/2018_Herzel_Ranging_Precision_FMCW_PLL.pdf} \\
\end{tabularx}
\textbf{Why it is here.} PLL phase noise plus receiver noise precision model.

\textbf{What to extract for this project.} Use for the precision/error-budget layer. Implement the impairment before the estimator whenever possible, then compare bias/variance/RMSE with the paper's analytical or measured behavior.

\textbf{Notebook prompt.} Write down the paper's assumptions in equation form. Mark each assumption as \emph{true in V0}, \emph{true only after hardware synchronization}, \emph{requires calibration}, or \emph{not true for two independent AWR boards}. This classification is more valuable than a generic summary.

\vspace{0.35em}\hrule\vspace{0.65em}

\clearpage
\section{P0 - phase noise / nonidealities}

\subsection*{15. Detailed Analysis and Modeling of Phase Noise and Systematic Phase Distortions in FMCW Radar Systems (2022)}
\begin{tabularx}{\textwidth}{@{}p{0.18\textwidth}X@{}}
\textbf{Priority/status} & P0 / \nolinkurl{included_restricted_user_supplied} \\
\textbf{Authors} & P. Tschapek et al. \\
\textbf{DOI / identifier} & \nolinkurl{10.1109/JMW.2022.3195574} \\
\textbf{Archive path} & \path{09_RESTRICTED_OR_CITATION_ONLY/user_supplied_restricted/2022_Tschapek_Detailed_Phase_Noise_Systematic_Distortions.pdf} \\
\end{tabularx}
\textbf{Why it is here.} Colored PN plus systematic phase-error Monte Carlo model.

\textbf{What to extract for this project.} Use for the precision/error-budget layer. Implement the impairment before the estimator whenever possible, then compare bias/variance/RMSE with the paper's analytical or measured behavior.

\textbf{Notebook prompt.} Write down the paper's assumptions in equation form. Mark each assumption as \emph{true in V0}, \emph{true only after hardware synchronization}, \emph{requires calibration}, or \emph{not true for two independent AWR boards}. This classification is more valuable than a generic summary.

\vspace{0.35em}\hrule\vspace{0.65em}

\clearpage
\section{P0 - precision / phase noise}

\subsection*{16. Performance Analysis of Cooperative FMCW Radar Distance Measurement Systems (2008)}
\begin{tabularx}{\textwidth}{@{}p{0.18\textwidth}X@{}}
\textbf{Priority/status} & P0 / \nolinkurl{citation_only} \\
\textbf{Authors} & S. Scheiblhofer et al. \\
\textbf{DOI / identifier} & \nolinkurl{10.1109/MWSYM.2008.4633118} \\
\textbf{Archive path} & not present \\
\end{tabularx}
\textbf{Why it is here.} Phase-noise performance comparison; web-accessible metadata, PDF not copied.

\textbf{What to extract for this project.} Use for the precision/error-budget layer. Implement the impairment before the estimator whenever possible, then compare bias/variance/RMSE with the paper's analytical or measured behavior.

\textbf{Notebook prompt.} Write down the paper's assumptions in equation form. Mark each assumption as \emph{true in V0}, \emph{true only after hardware synchronization}, \emph{requires calibration}, or \emph{not true for two independent AWR boards}. This classification is more valuable than a generic summary.

\vspace{0.35em}\hrule\vspace{0.65em}

\clearpage
\section{P1 - coded / multistatic synchronization}

\subsection*{17. On the Synchronization of Uncoupled Multistatic PMCW Radars (2024)}
\begin{tabularx}{\textwidth}{@{}p{0.18\textwidth}X@{}}
\textbf{Priority/status} & P1 / \nolinkurl{citation_only} \\
\textbf{Authors} & D. Werbunat et al. \\
\textbf{DOI / identifier} & \nolinkurl{10.1109/TMTT.2024.3359035} \\
\textbf{Archive path} & not present \\
\end{tabularx}
\textbf{Why it is here.} Digital carrier/timing recovery analog for future coded FMCW.

\textbf{What to extract for this project.} Use to formalize clock epoch offset, clock-rate/frequency offset, reciprocity, two-way cancellation, presynchronization, and calibration. Translate each distinct clock/path error into an explicit simulator parameter.

\textbf{Notebook prompt.} Write down the paper's assumptions in equation form. Mark each assumption as \emph{true in V0}, \emph{true only after hardware synchronization}, \emph{requires calibration}, or \emph{not true for two independent AWR boards}. This classification is more valuable than a generic summary.

\vspace{0.35em}\hrule\vspace{0.65em}

\clearpage
\section{P1 - coded FMCW}

\subsection*{18. Phase-Coded FMCW Automotive Radar: System Design and Interference Mitigation (2020)}
\begin{tabularx}{\textwidth}{@{}p{0.18\textwidth}X@{}}
\textbf{Priority/status} & P1 / \nolinkurl{included_acquired_research_copy} \\
\textbf{Authors} & F. Uysal \\
\textbf{DOI / identifier} & \nolinkurl{10.1109/TVT.2019.2953305} \\
\textbf{Archive path} & \path{04_CODED_FMCW_AND_MULTIRADAR/additional_acquired/2020_Uysal_Phase_Coded_FMCW_System_Design_Interference_Mitigation.pdf} \\
\end{tabularx}
\textbf{Why it is here.} Accepted manuscript known at TU Delft; not copied due repository download restriction.

\textbf{What to extract for this project.} Use only after the uncoded V0/V1 model is validated. Extract how waveform identity is embedded, how synchronization/code alignment is performed, and whether dechirp preserves the low-bandwidth receiver advantage.

\textbf{Notebook prompt.} Write down the paper's assumptions in equation form. Mark each assumption as \emph{true in V0}, \emph{true only after hardware synchronization}, \emph{requires calibration}, or \emph{not true for two independent AWR boards}. This classification is more valuable than a generic summary.

\vspace{0.35em}\hrule\vspace{0.65em}

\subsection*{19. System Level Synchronization of Phase-Coded FMCW Automotive Radars for RadCom (2020)}
\begin{tabularx}{\textwidth}{@{}p{0.18\textwidth}X@{}}
\textbf{Priority/status} & P1 / \nolinkurl{included_restricted_user_supplied} \\
\textbf{Authors} & F. Lampel et al. \\
\textbf{DOI / identifier} & \nolinkurl{10.23919/EuCAP48036.2020.9135417} \\
\textbf{Archive path} & \path{09_RESTRICTED_OR_CITATION_ONLY/user_supplied_restricted/2020_Lampel_System_Level_Synchronization_Phase_Coded_FMCW_RadCom.pdf} \\
\end{tabularx}
\textbf{Why it is here.} Phase-coded FMCW synchronization / RadCom.

\textbf{What to extract for this project.} Use to formalize clock epoch offset, clock-rate/frequency offset, reciprocity, two-way cancellation, presynchronization, and calibration. Translate each distinct clock/path error into an explicit simulator parameter.

\textbf{Notebook prompt.} Write down the paper's assumptions in equation form. Mark each assumption as \emph{true in V0}, \emph{true only after hardware synchronization}, \emph{requires calibration}, or \emph{not true for two independent AWR boards}. This classification is more valuable than a generic summary.

\vspace{0.35em}\hrule\vspace{0.65em}

\clearpage
\section{P1 - distributed coherent radar}

\subsection*{20. Coherent Signal Processing for Loosely Coupled Bistatic Radar (2021)}
\begin{tabularx}{\textwidth}{@{}p{0.18\textwidth}X@{}}
\textbf{Priority/status} & P1 / \nolinkurl{included_restricted_user_supplied} \\
\textbf{Authors} & M. Gottinger; P. Gulden; M. Vossiek \\
\textbf{DOI / identifier} & \nolinkurl{10.1109/TAES.2021.3050650} \\
\textbf{Archive path} & \path{09_RESTRICTED_OR_CITATION_ONLY/user_supplied_restricted/2021_Gottinger_Loosely_Coupled_Bistatic_Radar.pdf} \\
\end{tabularx}
\textbf{Why it is here.} Coherent processing despite only loose coupling.

\textbf{What to extract for this project.} Use to understand which observables survive separate oscillators, what degree of coherence is required, and how cross-path data should be retained/combined across nodes.

\textbf{Notebook prompt.} Write down the paper's assumptions in equation form. Mark each assumption as \emph{true in V0}, \emph{true only after hardware synchronization}, \emph{requires calibration}, or \emph{not true for two independent AWR boards}. This classification is more valuable than a generic summary.

\vspace{0.35em}\hrule\vspace{0.65em}

\clearpage
\section{P1 - radar network signal model}

\subsection*{21. Signal Model for Coherent Processing of Uncoupled and Low Frequency Coupled MIMO Radar Networks (2024)}
\begin{tabularx}{\textwidth}{@{}p{0.18\textwidth}X@{}}
\textbf{Priority/status} & P1 / \nolinkurl{citation_only} \\
\textbf{Authors} & V. Janoudi et al. \\
\textbf{DOI / identifier} & \nolinkurl{10.1109/JMW.2023.3334757} \\
\textbf{Archive path} & not present \\
\end{tabularx}
\textbf{Why it is here.} Strong modern nonideal signal model.

\textbf{What to extract for this project.} Use to map the abstract model to AWR2944/DCA1000 configuration fields, timing windows, sample rates, clocking, calibration, channel topology, and raw-data interpretation.

\textbf{Notebook prompt.} Write down the paper's assumptions in equation form. Mark each assumption as \emph{true in V0}, \emph{true only after hardware synchronization}, \emph{requires calibration}, or \emph{not true for two independent AWR boards}. This classification is more valuable than a generic summary.

\vspace{0.35em}\hrule\vspace{0.65em}

\clearpage
\section{P1 - radar network synchronization}

\subsection*{22. Over-the-Air Synchronization for Coherent Digital Automotive Radar Networks (2024)}
\begin{tabularx}{\textwidth}{@{}p{0.18\textwidth}X@{}}
\textbf{Priority/status} & P1 / \nolinkurl{citation_only} \\
\textbf{Authors} & L. Sigg et al. \\
\textbf{DOI / identifier} & \nolinkurl{10.1109/TRS.2024.3449333} \\
\textbf{Archive path} & not present \\
\end{tabularx}
\textbf{Why it is here.} CFO/SFO/timing correction.

\textbf{What to extract for this project.} Use to formalize clock epoch offset, clock-rate/frequency offset, reciprocity, two-way cancellation, presynchronization, and calibration. Translate each distinct clock/path error into an explicit simulator parameter.

\textbf{Notebook prompt.} Write down the paper's assumptions in equation form. Mark each assumption as \emph{true in V0}, \emph{true only after hardware synchronization}, \emph{requires calibration}, or \emph{not true for two independent AWR boards}. This classification is more valuable than a generic summary.

\vspace{0.35em}\hrule\vspace{0.65em}

\clearpage
\section{P1 - radar networks}

\subsection*{23. Coherent Automotive Radar Networks: The Next Generation of Radar-Based Imaging and Mapping (2021)}
\begin{tabularx}{\textwidth}{@{}p{0.18\textwidth}X@{}}
\textbf{Priority/status} & P1 / \nolinkurl{included_restricted_user_supplied} \\
\textbf{Authors} & M. Gottinger et al. \\
\textbf{DOI / identifier} & \nolinkurl{10.1109/JMW.2020.3034475} \\
\textbf{Archive path} & \path{09_RESTRICTED_OR_CITATION_ONLY/user_supplied_restricted/2021_Gottinger_Coherent_Automotive_Radar_Networks.pdf} \\
\end{tabularx}
\textbf{Why it is here.} Distributed radar classification and 76-77 GHz experiments.

\textbf{What to extract for this project.} Use to understand which observables survive separate oscillators, what degree of coherence is required, and how cross-path data should be retained/combined across nodes.

\textbf{Notebook prompt.} Write down the paper's assumptions in equation form. Mark each assumption as \emph{true in V0}, \emph{true only after hardware synchronization}, \emph{requires calibration}, or \emph{not true for two independent AWR boards}. This classification is more valuable than a generic summary.

\vspace{0.35em}\hrule\vspace{0.65em}

\clearpage
\section{P2 - additional acquired}

\subsection*{24. 2019 MathWorks Automotive Radar Development MATLAB (2019)}
\begin{tabularx}{\textwidth}{@{}p{0.18\textwidth}X@{}}
\textbf{Priority/status} & P2 / \nolinkurl{included_acquired_research_copy} \\
\textbf{Authors} & not recorded in compact catalog \\
\textbf{DOI / identifier} & not recorded \\
\textbf{Archive path} & \path{02_MATLAB_AND_SIMULATION/additional_acquired/2019_MathWorks_Automotive_Radar_Development_MATLAB.pdf} \\
\end{tabularx}
\textbf{Why it is here.} Acquired during the literature sweep and included as a research copy; see original source for licensing/redistribution terms.

\textbf{What to extract for this project.} Use to define simulator architecture, signal-domain fidelity, validation strategy, or efficient baseband implementation. Extract module boundaries, required state variables, and validation metrics.

\textbf{Notebook prompt.} Write down the paper's assumptions in equation form. Mark each assumption as \emph{true in V0}, \emph{true only after hardware synchronization}, \emph{requires calibration}, or \emph{not true for two independent AWR boards}. This classification is more valuable than a generic summary.

\vspace{0.35em}\hrule\vspace{0.65em}

\subsection*{25. 2022 Duerr Coherent Multistatic MIMO Radar Network Phase Noise Optimized Synthesis (2022)}
\begin{tabularx}{\textwidth}{@{}p{0.18\textwidth}X@{}}
\textbf{Priority/status} & P2 / \nolinkurl{included_acquired_research_copy} \\
\textbf{Authors} & not recorded in compact catalog \\
\textbf{DOI / identifier} & not recorded \\
\textbf{Archive path} & \path{04_CODED_FMCW_AND_MULTIRADAR/additional_acquired/2022_Duerr_Coherent_Multistatic_MIMO_Radar_Network_Phase_Noise_Optimized_Synthesis.pdf} \\
\end{tabularx}
\textbf{Why it is here.} Acquired during the literature sweep and included as a research copy; see original source for licensing/redistribution terms.

\textbf{What to extract for this project.} Use for the precision/error-budget layer. Implement the impairment before the estimator whenever possible, then compare bias/variance/RMSE with the paper's analytical or measured behavior.

\textbf{Notebook prompt.} Write down the paper's assumptions in equation form. Mark each assumption as \emph{true in V0}, \emph{true only after hardware synchronization}, \emph{requires calibration}, or \emph{not true for two independent AWR boards}. This classification is more valuable than a generic summary.

\vspace{0.35em}\hrule\vspace{0.65em}

\subsection*{26. 2024 Ahmadi et al Distributed Massive MIMO FMCW Radar Simulator Ray Tracing (2024)}
\begin{tabularx}{\textwidth}{@{}p{0.18\textwidth}X@{}}
\textbf{Priority/status} & P2 / \nolinkurl{included_acquired_research_copy} \\
\textbf{Authors} & not recorded in compact catalog \\
\textbf{DOI / identifier} & not recorded \\
\textbf{Archive path} & \path{02_MATLAB_AND_SIMULATION/additional_acquired/2024_Ahmadi_et_al_Distributed_Massive_MIMO_FMCW_Radar_Simulator_Ray_Tracing.pdf} \\
\end{tabularx}
\textbf{Why it is here.} Acquired during the literature sweep and included as a research copy; see original source for licensing/redistribution terms.

\textbf{What to extract for this project.} Use to define simulator architecture, signal-domain fidelity, validation strategy, or efficient baseband implementation. Extract module boundaries, required state variables, and validation metrics.

\textbf{Notebook prompt.} Write down the paper's assumptions in equation form. Mark each assumption as \emph{true in V0}, \emph{true only after hardware synchronization}, \emph{requires calibration}, or \emph{not true for two independent AWR boards}. This classification is more valuable than a generic summary.

\vspace{0.35em}\hrule\vspace{0.65em}

\clearpage
\section{P2 - extended bibliography}

\subsection*{27. 2014 Frischen Hasch Waldschmidt - Cooperative Radar: Uncorrelated Phase Noise (year n/a)}
\begin{tabularx}{\textwidth}{@{}p{0.18\textwidth}X@{}}
\textbf{Priority/status} & P2 / \nolinkurl{included_acquired_research_copy} \\
\textbf{Authors} & not recorded in compact catalog \\
\textbf{DOI / identifier} & not recorded \\
\textbf{Archive path} & \path{04_CODED_FMCW_AND_MULTIRADAR/additional_acquired/2014_Frischen_Hasch_Waldschmidt_Cooperative_Radar_Uncorrelated_Phase_Noise.pdf} \\
\end{tabularx}
\textbf{Why it is here.} Listed in prior acquisition staging; no separate PDF bytes available in this runtime. Physical research copy is included in the topical archive; check original publication license before redistribution.

\textbf{What to extract for this project.} Use for the precision/error-budget layer. Implement the impairment before the estimator whenever possible, then compare bias/variance/RMSE with the paper's analytical or measured behavior.

\textbf{Notebook prompt.} Write down the paper's assumptions in equation form. Mark each assumption as \emph{true in V0}, \emph{true only after hardware synchronization}, \emph{requires calibration}, or \emph{not true for two independent AWR boards}. This classification is more valuable than a generic summary.

\vspace{0.35em}\hrule\vspace{0.65em}

\subsection*{28. 2014 Rajan van der Veen - Joint Ranging and Synchronization for Anchorless Mobile Networks (year n/a)}
\begin{tabularx}{\textwidth}{@{}p{0.18\textwidth}X@{}}
\textbf{Priority/status} & P2 / \nolinkurl{included_acquired_research_copy} \\
\textbf{Authors} & not recorded in compact catalog \\
\textbf{DOI / identifier} & not recorded \\
\textbf{Archive path} & \path{03_TWO_WAY_TIME_TRANSFER_AND_SYNC/additional_acquired/2014_Rajan_van_der_Veen_Joint_Ranging_Synchronization_Anchorless_Mobile_Network.pdf} \\
\end{tabularx}
\textbf{Why it is here.} Listed in prior acquisition staging; no separate PDF bytes available in this runtime. Physical research copy is included in the topical archive; check original publication license before redistribution.

\textbf{What to extract for this project.} Use to formalize clock epoch offset, clock-rate/frequency offset, reciprocity, two-way cancellation, presynchronization, and calibration. Translate each distinct clock/path error into an explicit simulator parameter.

\textbf{Notebook prompt.} Write down the paper's assumptions in equation form. Mark each assumption as \emph{true in V0}, \emph{true only after hardware synchronization}, \emph{requires calibration}, or \emph{not true for two independent AWR boards}. This classification is more valuable than a generic summary.

\vspace{0.35em}\hrule\vspace{0.65em}

\subsection*{29. 2015 Dwivedi - Joint Ranging and Clock Parameter Estimation via Two-Way RTT (year n/a)}
\begin{tabularx}{\textwidth}{@{}p{0.18\textwidth}X@{}}
\textbf{Priority/status} & P2 / \nolinkurl{included_acquired_research_copy} \\
\textbf{Authors} & not recorded in compact catalog \\
\textbf{DOI / identifier} & not recorded \\
\textbf{Archive path} & \path{03_TWO_WAY_TIME_TRANSFER_AND_SYNC/additional_acquired/2015_Dwivedi_Joint_Ranging_Clock_Parameter_Estimation_Two_Way_RTT.pdf} \\
\end{tabularx}
\textbf{Why it is here.} Listed in prior acquisition staging; no separate PDF bytes available in this runtime. Physical research copy is included in the topical archive; check original publication license before redistribution.

\textbf{What to extract for this project.} Use to formalize clock epoch offset, clock-rate/frequency offset, reciprocity, two-way cancellation, presynchronization, and calibration. Translate each distinct clock/path error into an explicit simulator parameter.

\textbf{Notebook prompt.} Write down the paper's assumptions in equation form. Mark each assumption as \emph{true in V0}, \emph{true only after hardware synchronization}, \emph{requires calibration}, or \emph{not true for two independent AWR boards}. This classification is more valuable than a generic summary.

\vspace{0.35em}\hrule\vspace{0.65em}

\subsection*{30. 2015 ITU - Two-Way Satellite Time and Frequency Transfer (year n/a)}
\begin{tabularx}{\textwidth}{@{}p{0.18\textwidth}X@{}}
\textbf{Priority/status} & P2 / \nolinkurl{included_acquired_research_copy} \\
\textbf{Authors} & not recorded in compact catalog \\
\textbf{DOI / identifier} & not recorded \\
\textbf{Archive path} & \path{03_TWO_WAY_TIME_TRANSFER_AND_SYNC/additional_acquired/2015_ITU_TF1153_Two_Way_Satellite_Time_Frequency_Transfer.pdf} \\
\end{tabularx}
\textbf{Why it is here.} Listed in prior acquisition staging; no separate PDF bytes available in this runtime. Physical research copy is included in the topical archive; check original publication license before redistribution.

\textbf{What to extract for this project.} Use to formalize clock epoch offset, clock-rate/frequency offset, reciprocity, two-way cancellation, presynchronization, and calibration. Translate each distinct clock/path error into an explicit simulator parameter.

\textbf{Notebook prompt.} Write down the paper's assumptions in equation form. Mark each assumption as \emph{true in V0}, \emph{true only after hardware synchronization}, \emph{requires calibration}, or \emph{not true for two independent AWR boards}. This classification is more valuable than a generic summary.

\vspace{0.35em}\hrule\vspace{0.65em}

\subsection*{31. 2016 Lopez-Martinez et al. - Simulation of FMCW Radar with SDR (year n/a)}
\begin{tabularx}{\textwidth}{@{}p{0.18\textwidth}X@{}}
\textbf{Priority/status} & P2 / \nolinkurl{included_acquired_research_copy} \\
\textbf{Authors} & not recorded in compact catalog \\
\textbf{DOI / identifier} & not recorded \\
\textbf{Archive path} & \path{02_MATLAB_AND_SIMULATION/additional_acquired/2016_Lopez_Martinez_Vidal_Morera_Simulation_FMCW_Radar_SDR.pdf} \\
\end{tabularx}
\textbf{Why it is here.} Listed in prior acquisition staging; no separate PDF bytes available in this runtime. Physical research copy is included in the topical archive; check original publication license before redistribution.

\textbf{What to extract for this project.} Use to define simulator architecture, signal-domain fidelity, validation strategy, or efficient baseband implementation. Extract module boundaries, required state variables, and validation metrics.

\textbf{Notebook prompt.} Write down the paper's assumptions in equation form. Mark each assumption as \emph{true in V0}, \emph{true only after hardware synchronization}, \emph{requires calibration}, or \emph{not true for two independent AWR boards}. This classification is more valuable than a generic summary.

\vspace{0.35em}\hrule\vspace{0.65em}

\subsection*{32. 2016 Scheiblhofer - In-Chirp FSK Communication Between Cooperative 77-GHz Radar Stations (year n/a)}
\begin{tabularx}{\textwidth}{@{}p{0.18\textwidth}X@{}}
\textbf{Priority/status} & P2 / \nolinkurl{included_acquired_research_copy} \\
\textbf{Authors} & not recorded in compact catalog \\
\textbf{DOI / identifier} & not recorded \\
\textbf{Archive path} & \path{04_CODED_FMCW_AND_MULTIRADAR/additional_acquired/2016_Scheiblhofer_In_Chirp_FSK_Communication_Cooperative_77GHz_Radar.pdf} \\
\end{tabularx}
\textbf{Why it is here.} Listed in prior acquisition staging; no separate PDF bytes available in this runtime. Physical research copy is included in the topical archive; check original publication license before redistribution.

\textbf{What to extract for this project.} Use only after the uncoded V0/V1 model is validated. Extract how waveform identity is embedded, how synchronization/code alignment is performed, and whether dechirp preserves the low-bandwidth receiver advantage.

\textbf{Notebook prompt.} Write down the paper's assumptions in equation form. Mark each assumption as \emph{true in V0}, \emph{true only after hardware synchronization}, \emph{requires calibration}, or \emph{not true for two independent AWR boards}. This classification is more valuable than a generic summary.

\vspace{0.35em}\hrule\vspace{0.65em}

\subsection*{33. 2017 TI - Complex Baseband Architecture in FMCW Radar Systems (year n/a)}
\begin{tabularx}{\textwidth}{@{}p{0.18\textwidth}X@{}}
\textbf{Priority/status} & P2 / \nolinkurl{included_acquired_research_copy} \\
\textbf{Authors} & not recorded in compact catalog \\
\textbf{DOI / identifier} & not recorded \\
\textbf{Archive path} & \path{01_FMCW_FUNDAMENTALS/classic_and_signal_processing/2017_TI_Complex_Baseband_Architecture_in_FMCW_Radar_Systems.pdf} \\
\end{tabularx}
\textbf{Why it is here.} Listed in prior acquisition staging; no separate PDF bytes available in this runtime. Physical research copy is included in the topical archive; check original publication license before redistribution.

\textbf{What to extract for this project.} Reference for FMCW/radar fundamentals or adjacent methods. Use it to sanity-check notation, signal relationships, and assumptions when its topic becomes active in the implementation.

\textbf{Notebook prompt.} Write down the paper's assumptions in equation form. Mark each assumption as \emph{true in V0}, \emph{true only after hardware synchronization}, \emph{requires calibration}, or \emph{not true for two independent AWR boards}. This classification is more valuable than a generic summary.

\vspace{0.35em}\hrule\vspace{0.65em}

\subsection*{34. 2018 Kazaz - Joint Ranging and Clock Synchronization for Dense IoT (year n/a)}
\begin{tabularx}{\textwidth}{@{}p{0.18\textwidth}X@{}}
\textbf{Priority/status} & P2 / \nolinkurl{included_acquired_research_copy} \\
\textbf{Authors} & not recorded in compact catalog \\
\textbf{DOI / identifier} & not recorded \\
\textbf{Archive path} & \path{03_TWO_WAY_TIME_TRANSFER_AND_SYNC/additional_acquired/2018_Kazaz_Joint_Ranging_Clock_Synchronization_Dense_IoT.pdf} \\
\end{tabularx}
\textbf{Why it is here.} Listed in prior acquisition staging; no separate PDF bytes available in this runtime. Physical research copy is included in the topical archive; check original publication license before redistribution.

\textbf{What to extract for this project.} Use to formalize clock epoch offset, clock-rate/frequency offset, reciprocity, two-way cancellation, presynchronization, and calibration. Translate each distinct clock/path error into an explicit simulator parameter.

\textbf{Notebook prompt.} Write down the paper's assumptions in equation form. Mark each assumption as \emph{true in V0}, \emph{true only after hardware synchronization}, \emph{requires calibration}, or \emph{not true for two independent AWR boards}. This classification is more valuable than a generic summary.

\vspace{0.35em}\hrule\vspace{0.65em}

\subsection*{35. 2018 Park et al. - Leakage Mitigation / Internal Delay Compensation in FMCW (year n/a)}
\begin{tabularx}{\textwidth}{@{}p{0.18\textwidth}X@{}}
\textbf{Priority/status} & P2 / \nolinkurl{included_acquired_research_copy} \\
\textbf{Authors} & not recorded in compact catalog \\
\textbf{DOI / identifier} & not recorded \\
\textbf{Archive path} & \path{05_PRECISION_LIMITS_AND_NONIDEALITIES/additional_acquired/2018_Park_Park_Park_Leakage_Mitigation_Internal_Delay_Compensation_FMCW.pdf} \\
\end{tabularx}
\textbf{Why it is here.} Listed in prior acquisition staging; no separate PDF bytes available in this runtime. Physical research copy is included in the topical archive; check original publication license before redistribution.

\textbf{What to extract for this project.} Use for the precision/error-budget layer. Implement the impairment before the estimator whenever possible, then compare bias/variance/RMSE with the paper's analytical or measured behavior.

\textbf{Notebook prompt.} Write down the paper's assumptions in equation form. Mark each assumption as \emph{true in V0}, \emph{true only after hardware synchronization}, \emph{requires calibration}, or \emph{not true for two independent AWR boards}. This classification is more valuable than a generic summary.

\vspace{0.35em}\hrule\vspace{0.65em}

\subsection*{36. 2018 Svensson - High Resolution Frequency Estimation for FMCW Radar (thesis) (year n/a)}
\begin{tabularx}{\textwidth}{@{}p{0.18\textwidth}X@{}}
\textbf{Priority/status} & P2 / \nolinkurl{included_acquired_research_copy} \\
\textbf{Authors} & not recorded in compact catalog \\
\textbf{DOI / identifier} & not recorded \\
\textbf{Archive path} & \path{05_PRECISION_LIMITS_AND_NONIDEALITIES/additional_acquired/2018_Svensson_High_Resolution_Frequency_Estimation_FMCW_Radar_Thesis.pdf} \\
\end{tabularx}
\textbf{Why it is here.} Listed in prior acquisition staging; no separate PDF bytes available in this runtime. Physical research copy is included in the topical archive; check original publication license before redistribution.

\textbf{What to extract for this project.} Use for the precision/error-budget layer. Implement the impairment before the estimator whenever possible, then compare bias/variance/RMSE with the paper's analytical or measured behavior.

\textbf{Notebook prompt.} Write down the paper's assumptions in equation form. Mark each assumption as \emph{true in V0}, \emph{true only after hardware synchronization}, \emph{requires calibration}, or \emph{not true for two independent AWR boards}. This classification is more valuable than a generic summary.

\vspace{0.35em}\hrule\vspace{0.65em}

\subsection*{37. 2019 Mueller Diewald - Cooperative Radar Signature Method for Unambiguity (year n/a)}
\begin{tabularx}{\textwidth}{@{}p{0.18\textwidth}X@{}}
\textbf{Priority/status} & P2 / \nolinkurl{included_acquired_research_copy} \\
\textbf{Authors} & not recorded in compact catalog \\
\textbf{DOI / identifier} & not recorded \\
\textbf{Archive path} & \path{04_CODED_FMCW_AND_MULTIRADAR/additional_acquired/2019_Mueller_Diewald_Cooperative_Radar_Signature_Unambiguity.pdf} \\
\end{tabularx}
\textbf{Why it is here.} Listed in prior acquisition staging; no separate PDF bytes available in this runtime. Physical research copy is included in the topical archive; check original publication license before redistribution.

\textbf{What to extract for this project.} Use only after the uncoded V0/V1 model is validated. Extract how waveform identity is embedded, how synchronization/code alignment is performed, and whether dechirp preserves the low-bandwidth receiver advantage.

\textbf{Notebook prompt.} Write down the paper's assumptions in equation form. Mark each assumption as \emph{true in V0}, \emph{true only after hardware synchronization}, \emph{requires calibration}, or \emph{not true for two independent AWR boards}. This classification is more valuable than a generic summary.

\vspace{0.35em}\hrule\vspace{0.65em}

\subsection*{38. 2019 Wang et al. - Range Accuracy of FMCW Radar with Source Nonlinearity (year n/a)}
\begin{tabularx}{\textwidth}{@{}p{0.18\textwidth}X@{}}
\textbf{Priority/status} & P2 / \nolinkurl{included_acquired_research_copy} \\
\textbf{Authors} & not recorded in compact catalog \\
\textbf{DOI / identifier} & not recorded \\
\textbf{Archive path} & \path{05_PRECISION_LIMITS_AND_NONIDEALITIES/additional_acquired/2019_Wang_et_al_Range_Accuracy_FMCW_Source_Nonlinearity.pdf} \\
\end{tabularx}
\textbf{Why it is here.} Listed in prior acquisition staging; no separate PDF bytes available in this runtime. Physical research copy is included in the topical archive; check original publication license before redistribution.

\textbf{What to extract for this project.} Use for the precision/error-budget layer. Implement the impairment before the estimator whenever possible, then compare bias/variance/RMSE with the paper's analytical or measured behavior.

\textbf{Notebook prompt.} Write down the paper's assumptions in equation form. Mark each assumption as \emph{true in V0}, \emph{true only after hardware synchronization}, \emph{requires calibration}, or \emph{not true for two independent AWR boards}. This classification is more valuable than a generic summary.

\vspace{0.35em}\hrule\vspace{0.65em}

\subsection*{39. 2020 Infineon - Phase Coded FMCW Radar patent (year n/a)}
\begin{tabularx}{\textwidth}{@{}p{0.18\textwidth}X@{}}
\textbf{Priority/status} & P2 / \nolinkurl{included_acquired_research_copy} \\
\textbf{Authors} & not recorded in compact catalog \\
\textbf{DOI / identifier} & not recorded \\
\textbf{Archive path} & \path{04_CODED_FMCW_AND_MULTIRADAR/additional_acquired/2020_Infineon_US20200132825A1_Phase_Coded_FMCW_Radar.pdf} \\
\end{tabularx}
\textbf{Why it is here.} Listed in prior acquisition staging; no separate PDF bytes available in this runtime. Physical research copy is included in the topical archive; check original publication license before redistribution.

\textbf{What to extract for this project.} Use only after the uncoded V0/V1 model is validated. Extract how waveform identity is embedded, how synchronization/code alignment is performed, and whether dechirp preserves the low-bandwidth receiver advantage.

\textbf{Notebook prompt.} Write down the paper's assumptions in equation form. Mark each assumption as \emph{true in V0}, \emph{true only after hardware synchronization}, \emph{requires calibration}, or \emph{not true for two independent AWR boards}. This classification is more valuable than a generic summary.

\vspace{0.35em}\hrule\vspace{0.65em}

\subsection*{40. 2020 UIUC - FMCW Radar lecture (year n/a)}
\begin{tabularx}{\textwidth}{@{}p{0.18\textwidth}X@{}}
\textbf{Priority/status} & P2 / \nolinkurl{included_acquired_research_copy} \\
\textbf{Authors} & not recorded in compact catalog \\
\textbf{DOI / identifier} & not recorded \\
\textbf{Archive path} & \path{02_MATLAB_AND_SIMULATION/additional_acquired/2020_UIUC_FMCW_Radar_Lecture.pdf} \\
\end{tabularx}
\textbf{Why it is here.} Listed in prior acquisition staging; no separate PDF bytes available in this runtime. Physical research copy is included in the topical archive; check original publication license before redistribution.

\textbf{What to extract for this project.} Reference for FMCW/radar fundamentals or adjacent methods. Use it to sanity-check notation, signal relationships, and assumptions when its topic becomes active in the implementation.

\textbf{Notebook prompt.} Write down the paper's assumptions in equation form. Mark each assumption as \emph{true in V0}, \emph{true only after hardware synchronization}, \emph{requires calibration}, or \emph{not true for two independent AWR boards}. This classification is more valuable than a generic summary.

\vspace{0.35em}\hrule\vspace{0.65em}

\subsection*{41. 2020 Uysal Orru - Phase-Coded FMCW Automotive Radar: Application and Challenges (year n/a)}
\begin{tabularx}{\textwidth}{@{}p{0.18\textwidth}X@{}}
\textbf{Priority/status} & P2 / \nolinkurl{included_acquired_research_copy} \\
\textbf{Authors} & not recorded in compact catalog \\
\textbf{DOI / identifier} & not recorded \\
\textbf{Archive path} & \path{04_CODED_FMCW_AND_MULTIRADAR/additional_acquired/2020_Uysal_Orru_Phase_Coded_FMCW_Application_and_Challenges.pdf} \\
\end{tabularx}
\textbf{Why it is here.} Listed in prior acquisition staging; no separate PDF bytes available in this runtime. Physical research copy is included in the topical archive; check original publication license before redistribution.

\textbf{What to extract for this project.} Use only after the uncoded V0/V1 model is validated. Extract how waveform identity is embedded, how synchronization/code alignment is performed, and whether dechirp preserves the low-bandwidth receiver advantage.

\textbf{Notebook prompt.} Write down the paper's assumptions in equation form. Mark each assumption as \emph{true in V0}, \emph{true only after hardware synchronization}, \emph{requires calibration}, or \emph{not true for two independent AWR boards}. This classification is more valuable than a generic summary.

\vspace{0.35em}\hrule\vspace{0.65em}

\subsection*{42. 2022 Architectures and Synchronization Techniques for Distributed Satellite Systems - Survey (year n/a)}
\begin{tabularx}{\textwidth}{@{}p{0.18\textwidth}X@{}}
\textbf{Priority/status} & P2 / \nolinkurl{included_acquired_research_copy} \\
\textbf{Authors} & not recorded in compact catalog \\
\textbf{DOI / identifier} & not recorded \\
\textbf{Archive path} & \path{03_TWO_WAY_TIME_TRANSFER_AND_SYNC/additional_acquired/2022_Architectures_and_Synchronization_Techniques_Distributed_Satellite_Systems_Survey.pdf} \\
\end{tabularx}
\textbf{Why it is here.} Listed in prior acquisition staging; no separate PDF bytes available in this runtime. Physical research copy is included in the topical archive; check original publication license before redistribution.

\textbf{What to extract for this project.} Use to formalize clock epoch offset, clock-rate/frequency offset, reciprocity, two-way cancellation, presynchronization, and calibration. Translate each distinct clock/path error into an explicit simulator parameter.

\textbf{Notebook prompt.} Write down the paper's assumptions in equation form. Mark each assumption as \emph{true in V0}, \emph{true only after hardware synchronization}, \emph{requires calibration}, or \emph{not true for two independent AWR boards}. This classification is more valuable than a generic summary.

\vspace{0.35em}\hrule\vspace{0.65em}

\subsection*{43. 2022 Duerr - Coherent Multistatic MIMO Radar Network with Phase-Noise-Optimized Synthesis (year n/a)}
\begin{tabularx}{\textwidth}{@{}p{0.18\textwidth}X@{}}
\textbf{Priority/status} & P2 / \nolinkurl{included_acquired_research_copy} \\
\textbf{Authors} & not recorded in compact catalog \\
\textbf{DOI / identifier} & not recorded \\
\textbf{Archive path} & \path{06_TI_AWR2944_IMPLEMENTATION/chirp_config_calibration_and_MIMO/TI_MIMO_Radar.pdf} \\
\end{tabularx}
\textbf{Why it is here.} Listed in prior acquisition staging; no separate PDF bytes available in this runtime. Physical research copy is included in the topical archive; check original publication license before redistribution.

\textbf{What to extract for this project.} Use for the precision/error-budget layer. Implement the impairment before the estimator whenever possible, then compare bias/variance/RMSE with the paper's analytical or measured behavior.

\textbf{Notebook prompt.} Write down the paper's assumptions in equation form. Mark each assumption as \emph{true in V0}, \emph{true only after hardware synchronization}, \emph{requires calibration}, or \emph{not true for two independent AWR boards}. This classification is more valuable than a generic summary.

\vspace{0.35em}\hrule\vspace{0.65em}

\subsection*{44. 2022 MIT - Radar and Coherent Sensor Processing lecture (year n/a)}
\begin{tabularx}{\textwidth}{@{}p{0.18\textwidth}X@{}}
\textbf{Priority/status} & P2 / \nolinkurl{included_acquired_research_copy} \\
\textbf{Authors} & not recorded in compact catalog \\
\textbf{DOI / identifier} & not recorded \\
\textbf{Archive path} & \path{02_MATLAB_AND_SIMULATION/additional_acquired/2022_MIT_Radar_and_Coherent_Sensor_Processing_Lecture.pdf} \\
\end{tabularx}
\textbf{Why it is here.} Listed in prior acquisition staging; no separate PDF bytes available in this runtime. Physical research copy is included in the topical archive; check original publication license before redistribution.

\textbf{What to extract for this project.} Use to understand which observables survive separate oscillators, what degree of coherence is required, and how cross-path data should be retained/combined across nodes.

\textbf{Notebook prompt.} Write down the paper's assumptions in equation form. Mark each assumption as \emph{true in V0}, \emph{true only after hardware synchronization}, \emph{requires calibration}, or \emph{not true for two independent AWR boards}. This classification is more valuable than a generic summary.

\vspace{0.35em}\hrule\vspace{0.65em}

\subsection*{45. 2022 TI - Interference Mitigation for AWR/IWR Devices (year n/a)}
\begin{tabularx}{\textwidth}{@{}p{0.18\textwidth}X@{}}
\textbf{Priority/status} & P2 / \nolinkurl{included_acquired_research_copy} \\
\textbf{Authors} & not recorded in compact catalog \\
\textbf{DOI / identifier} & not recorded \\
\textbf{Archive path} & \path{01_FMCW_FUNDAMENTALS/classic_and_signal_processing/2022_TI_Interference_Mitigation_for_AWR_IWR_Devices.pdf} \\
\end{tabularx}
\textbf{Why it is here.} Listed in prior acquisition staging; no separate PDF bytes available in this runtime. Physical research copy is included in the topical archive; check original publication license before redistribution.

\textbf{What to extract for this project.} Reference for FMCW/radar fundamentals or adjacent methods. Use it to sanity-check notation, signal relationships, and assumptions when its topic becomes active in the implementation.

\textbf{Notebook prompt.} Write down the paper's assumptions in equation form. Mark each assumption as \emph{true in V0}, \emph{true only after hardware synchronization}, \emph{requires calibration}, or \emph{not true for two independent AWR boards}. This classification is more valuable than a generic summary.

\vspace{0.35em}\hrule\vspace{0.65em}

\subsection*{46. 2023 Tagliaferri et al. - Cooperative Coherent Multistatic Imaging and Phase Synchronization (year n/a)}
\begin{tabularx}{\textwidth}{@{}p{0.18\textwidth}X@{}}
\textbf{Priority/status} & P2 / \nolinkurl{included_acquired_research_copy} \\
\textbf{Authors} & not recorded in compact catalog \\
\textbf{DOI / identifier} & not recorded \\
\textbf{Archive path} & \path{03_TWO_WAY_TIME_TRANSFER_AND_SYNC/additional_acquired/2023_Tagliaferri_et_al_Cooperative_Coherent_Multistatic_Imaging_Phase_Synchronization.pdf} \\
\end{tabularx}
\textbf{Why it is here.} Listed in prior acquisition staging; no separate PDF bytes available in this runtime. Physical research copy is included in the topical archive; check original publication license before redistribution.

\textbf{What to extract for this project.} Use to formalize clock epoch offset, clock-rate/frequency offset, reciprocity, two-way cancellation, presynchronization, and calibration. Translate each distinct clock/path error into an explicit simulator parameter.

\textbf{Notebook prompt.} Write down the paper's assumptions in equation form. Mark each assumption as \emph{true in V0}, \emph{true only after hardware synchronization}, \emph{requires calibration}, or \emph{not true for two independent AWR boards}. This classification is more valuable than a generic summary.

\vspace{0.35em}\hrule\vspace{0.65em}

\subsection*{47. 2024 Ahmadi et al. - Distributed Massive MIMO FMCW Radar Simulator with Ray Tracing (year n/a)}
\begin{tabularx}{\textwidth}{@{}p{0.18\textwidth}X@{}}
\textbf{Priority/status} & P2 / \nolinkurl{included_acquired_research_copy} \\
\textbf{Authors} & not recorded in compact catalog \\
\textbf{DOI / identifier} & not recorded \\
\textbf{Archive path} & \path{06_TI_AWR2944_IMPLEMENTATION/chirp_config_calibration_and_MIMO/TI_MIMO_Radar.pdf} \\
\end{tabularx}
\textbf{Why it is here.} Listed in prior acquisition staging; no separate PDF bytes available in this runtime. Physical research copy is included in the topical archive; check original publication license before redistribution.

\textbf{What to extract for this project.} Use to define simulator architecture, signal-domain fidelity, validation strategy, or efficient baseband implementation. Extract module boundaries, required state variables, and validation metrics.

\textbf{Notebook prompt.} Write down the paper's assumptions in equation form. Mark each assumption as \emph{true in V0}, \emph{true only after hardware synchronization}, \emph{requires calibration}, or \emph{not true for two independent AWR boards}. This classification is more valuable than a generic summary.

\vspace{0.35em}\hrule\vspace{0.65em}

\subsection*{48. 2024 Chen Niknejad - RF-Domain Leakage Cancellation for FMCW Radars (year n/a)}
\begin{tabularx}{\textwidth}{@{}p{0.18\textwidth}X@{}}
\textbf{Priority/status} & P2 / \nolinkurl{included_acquired_research_copy} \\
\textbf{Authors} & not recorded in compact catalog \\
\textbf{DOI / identifier} & not recorded \\
\textbf{Archive path} & \path{05_PRECISION_LIMITS_AND_NONIDEALITIES/additional_acquired/2024_Chen_Niknejad_RF_Domain_Leakage_Cancellation_FMCW_Radars.pdf} \\
\end{tabularx}
\textbf{Why it is here.} Listed in prior acquisition staging; no separate PDF bytes available in this runtime. Physical research copy is included in the topical archive; check original publication license before redistribution.

\textbf{What to extract for this project.} Use for the precision/error-budget layer. Implement the impairment before the estimator whenever possible, then compare bias/variance/RMSE with the paper's analytical or measured behavior.

\textbf{Notebook prompt.} Write down the paper's assumptions in equation form. Mark each assumption as \emph{true in V0}, \emph{true only after hardware synchronization}, \emph{requires calibration}, or \emph{not true for two independent AWR boards}. This classification is more valuable than a generic summary.

\vspace{0.35em}\hrule\vspace{0.65em}

\subsection*{49. 2024 Huang et al. - Overview of Millimeter-Wave Radar Modeling Methods (year n/a)}
\begin{tabularx}{\textwidth}{@{}p{0.18\textwidth}X@{}}
\textbf{Priority/status} & P2 / \nolinkurl{included_acquired_research_copy} \\
\textbf{Authors} & not recorded in compact catalog \\
\textbf{DOI / identifier} & not recorded \\
\textbf{Archive path} & \path{02_MATLAB_AND_SIMULATION/additional_acquired/2024_Huang_et_al_Overview_Millimeter_Wave_Radar_Modeling_Methods.pdf} \\
\end{tabularx}
\textbf{Why it is here.} Listed in prior acquisition staging; no separate PDF bytes available in this runtime. Physical research copy is included in the topical archive; check original publication license before redistribution.

\textbf{What to extract for this project.} Use to define simulator architecture, signal-domain fidelity, validation strategy, or efficient baseband implementation. Extract module boundaries, required state variables, and validation metrics.

\textbf{Notebook prompt.} Write down the paper's assumptions in equation form. Mark each assumption as \emph{true in V0}, \emph{true only after hardware synchronization}, \emph{requires calibration}, or \emph{not true for two independent AWR boards}. This classification is more valuable than a generic summary.

\vspace{0.35em}\hrule\vspace{0.65em}

\subsection*{50. 2024 Kou et al. - Distributed PMCW Radar Network Phase Noise (year n/a)}
\begin{tabularx}{\textwidth}{@{}p{0.18\textwidth}X@{}}
\textbf{Priority/status} & P2 / \nolinkurl{included_acquired_research_copy} \\
\textbf{Authors} & not recorded in compact catalog \\
\textbf{DOI / identifier} & not recorded \\
\textbf{Archive path} & \path{04_CODED_FMCW_AND_MULTIRADAR/additional_acquired/2024_Kou_Bauduin_Bourdoux_Pollin_Distributed_PMCW_Radar_Network_Phase_Noise.pdf} \\
\end{tabularx}
\textbf{Why it is here.} Listed in prior acquisition staging; no separate PDF bytes available in this runtime. Physical research copy is included in the topical archive; check original publication license before redistribution.

\textbf{What to extract for this project.} Use only after the uncoded V0/V1 model is validated. Extract how waveform identity is embedded, how synchronization/code alignment is performed, and whether dechirp preserves the low-bandwidth receiver advantage.

\textbf{Notebook prompt.} Write down the paper's assumptions in equation form. Mark each assumption as \emph{true in V0}, \emph{true only after hardware synchronization}, \emph{requires calibration}, or \emph{not true for two independent AWR boards}. This classification is more valuable than a generic summary.

\vspace{0.35em}\hrule\vspace{0.65em}

\subsection*{51. 2024 Multi-Objective Distributed Beamforming with High-Accuracy Synchronization (year n/a)}
\begin{tabularx}{\textwidth}{@{}p{0.18\textwidth}X@{}}
\textbf{Priority/status} & P2 / \nolinkurl{included_acquired_research_copy} \\
\textbf{Authors} & not recorded in compact catalog \\
\textbf{DOI / identifier} & not recorded \\
\textbf{Archive path} & \path{03_TWO_WAY_TIME_TRANSFER_AND_SYNC/additional_acquired/2024_Multi_Objective_Distributed_Beamforming_High_Accuracy_Sync.pdf} \\
\end{tabularx}
\textbf{Why it is here.} Listed in prior acquisition staging; no separate PDF bytes available in this runtime. Physical research copy is included in the topical archive; check original publication license before redistribution.

\textbf{What to extract for this project.} Use to formalize clock epoch offset, clock-rate/frequency offset, reciprocity, two-way cancellation, presynchronization, and calibration. Translate each distinct clock/path error into an explicit simulator parameter.

\textbf{Notebook prompt.} Write down the paper's assumptions in equation form. Mark each assumption as \emph{true in V0}, \emph{true only after hardware synchronization}, \emph{requires calibration}, or \emph{not true for two independent AWR boards}. This classification is more valuable than a generic summary.

\vspace{0.35em}\hrule\vspace{0.65em}

\subsection*{52. 2025 BIPM - Time Dissemination/TWSTFT lecture (year n/a)}
\begin{tabularx}{\textwidth}{@{}p{0.18\textwidth}X@{}}
\textbf{Priority/status} & P2 / \nolinkurl{included_acquired_research_copy} \\
\textbf{Authors} & not recorded in compact catalog \\
\textbf{DOI / identifier} & not recorded \\
\textbf{Archive path} & \path{03_TWO_WAY_TIME_TRANSFER_AND_SYNC/additional_acquired/2025_BIPM_Time_Dissemination_TWSTFT_Lecture.pdf} \\
\end{tabularx}
\textbf{Why it is here.} Listed in prior acquisition staging; no separate PDF bytes available in this runtime. Physical research copy is included in the topical archive; check original publication license before redistribution.

\textbf{What to extract for this project.} Use to formalize clock epoch offset, clock-rate/frequency offset, reciprocity, two-way cancellation, presynchronization, and calibration. Translate each distinct clock/path error into an explicit simulator parameter.

\textbf{Notebook prompt.} Write down the paper's assumptions in equation form. Mark each assumption as \emph{true in V0}, \emph{true only after hardware synchronization}, \emph{requires calibration}, or \emph{not true for two independent AWR boards}. This classification is more valuable than a generic summary.

\vspace{0.35em}\hrule\vspace{0.65em}

\subsection*{53. 2025 DLR - Investigation of CW/LFM Waveforms for Bi/Multistatic Radar Synchronization (year n/a)}
\begin{tabularx}{\textwidth}{@{}p{0.18\textwidth}X@{}}
\textbf{Priority/status} & P2 / \nolinkurl{included_acquired_research_copy} \\
\textbf{Authors} & not recorded in compact catalog \\
\textbf{DOI / identifier} & not recorded \\
\textbf{Archive path} & \path{03_TWO_WAY_TIME_TRANSFER_AND_SYNC/additional_acquired/2025_DLR_Investigation_CW_LFM_Waveforms_Bi_Multistatic_Radar_Synchronization.pdf} \\
\end{tabularx}
\textbf{Why it is here.} Listed in prior acquisition staging; no separate PDF bytes available in this runtime. Physical research copy is included in the topical archive; check original publication license before redistribution.

\textbf{What to extract for this project.} Use to formalize clock epoch offset, clock-rate/frequency offset, reciprocity, two-way cancellation, presynchronization, and calibration. Translate each distinct clock/path error into an explicit simulator parameter.

\textbf{Notebook prompt.} Write down the paper's assumptions in equation form. Mark each assumption as \emph{true in V0}, \emph{true only after hardware synchronization}, \emph{requires calibration}, or \emph{not true for two independent AWR boards}. This classification is more valuable than a generic summary.

\vspace{0.35em}\hrule\vspace{0.65em}

\subsection*{54. 2025 Eid - Novel Simulation Technique for UWB FMCW Radar (year n/a)}
\begin{tabularx}{\textwidth}{@{}p{0.18\textwidth}X@{}}
\textbf{Priority/status} & P2 / \nolinkurl{included_acquired_research_copy} \\
\textbf{Authors} & not recorded in compact catalog \\
\textbf{DOI / identifier} & not recorded \\
\textbf{Archive path} & \path{02_MATLAB_AND_SIMULATION/additional_acquired/2025_Eid_Novel_Simulation_Technique_UWB_FMCW_Radar.pdf} \\
\end{tabularx}
\textbf{Why it is here.} Listed in prior acquisition staging; no separate PDF bytes available in this runtime. Physical research copy is included in the topical archive; check original publication license before redistribution.

\textbf{What to extract for this project.} Use to define simulator architecture, signal-domain fidelity, validation strategy, or efficient baseband implementation. Extract module boundaries, required state variables, and validation metrics.

\textbf{Notebook prompt.} Write down the paper's assumptions in equation form. Mark each assumption as \emph{true in V0}, \emph{true only after hardware synchronization}, \emph{requires calibration}, or \emph{not true for two independent AWR boards}. This classification is more valuable than a generic summary.

\vspace{0.35em}\hrule\vspace{0.65em}

\subsection*{55. 2025 Han Meng Masouros - OTA Time/Frequency Synchronization for Distributed ISAC (year n/a)}
\begin{tabularx}{\textwidth}{@{}p{0.18\textwidth}X@{}}
\textbf{Priority/status} & P2 / \nolinkurl{included_acquired_research_copy} \\
\textbf{Authors} & not recorded in compact catalog \\
\textbf{DOI / identifier} & not recorded \\
\textbf{Archive path} & \path{03_TWO_WAY_TIME_TRANSFER_AND_SYNC/additional_acquired/2025_Han_Meng_Masouros_OTA_Time_Frequency_Synchronization_Distributed_ISAC.pdf} \\
\end{tabularx}
\textbf{Why it is here.} Listed in prior acquisition staging; no separate PDF bytes available in this runtime. Physical research copy is included in the topical archive; check original publication license before redistribution.

\textbf{What to extract for this project.} Use to formalize clock epoch offset, clock-rate/frequency offset, reciprocity, two-way cancellation, presynchronization, and calibration. Translate each distinct clock/path error into an explicit simulator parameter.

\textbf{Notebook prompt.} Write down the paper's assumptions in equation form. Mark each assumption as \emph{true in V0}, \emph{true only after hardware synchronization}, \emph{requires calibration}, or \emph{not true for two independent AWR boards}. This classification is more valuable than a generic summary.

\vspace{0.35em}\hrule\vspace{0.65em}

\subsection*{56. 2025 Multistatic Radar Performance with Distributed Wireless Synchronization / BCRLB (year n/a)}
\begin{tabularx}{\textwidth}{@{}p{0.18\textwidth}X@{}}
\textbf{Priority/status} & P2 / \nolinkurl{included_acquired_research_copy} \\
\textbf{Authors} & not recorded in compact catalog \\
\textbf{DOI / identifier} & not recorded \\
\textbf{Archive path} & \path{03_TWO_WAY_TIME_TRANSFER_AND_SYNC/additional_acquired/2025_Multistatic_Radar_Performance_Distributed_Wireless_Synchronization_BCRLB.pdf} \\
\end{tabularx}
\textbf{Why it is here.} Listed in prior acquisition staging; no separate PDF bytes available in this runtime. Physical research copy is included in the topical archive; check original publication license before redistribution.

\textbf{What to extract for this project.} Use to formalize clock epoch offset, clock-rate/frequency offset, reciprocity, two-way cancellation, presynchronization, and calibration. Translate each distinct clock/path error into an explicit simulator parameter.

\textbf{Notebook prompt.} Write down the paper's assumptions in equation form. Mark each assumption as \emph{true in V0}, \emph{true only after hardware synchronization}, \emph{requires calibration}, or \emph{not true for two independent AWR boards}. This classification is more valuable than a generic summary.

\vspace{0.35em}\hrule\vspace{0.65em}

\subsection*{57. 2025 Real-Time High-Accuracy Digital Wireless Time/Frequency/Phase Synchronization (year n/a)}
\begin{tabularx}{\textwidth}{@{}p{0.18\textwidth}X@{}}
\textbf{Priority/status} & P2 / \nolinkurl{included_acquired_research_copy} \\
\textbf{Authors} & not recorded in compact catalog \\
\textbf{DOI / identifier} & not recorded \\
\textbf{Archive path} & \path{03_TWO_WAY_TIME_TRANSFER_AND_SYNC/additional_acquired/2025_Real_Time_High_Accuracy_Digital_Wireless_Time_Frequency_Phase_Synchronization.pdf} \\
\end{tabularx}
\textbf{Why it is here.} Listed in prior acquisition staging; no separate PDF bytes available in this runtime. Physical research copy is included in the topical archive; check original publication license before redistribution.

\textbf{What to extract for this project.} Use to formalize clock epoch offset, clock-rate/frequency offset, reciprocity, two-way cancellation, presynchronization, and calibration. Translate each distinct clock/path error into an explicit simulator parameter.

\textbf{Notebook prompt.} Write down the paper's assumptions in equation form. Mark each assumption as \emph{true in V0}, \emph{true only after hardware synchronization}, \emph{requires calibration}, or \emph{not true for two independent AWR boards}. This classification is more valuable than a generic summary.

\vspace{0.35em}\hrule\vspace{0.65em}

\subsection*{58. 2026 UCSB - Introduction to mmWave Radar Sensing Lab Notes (year n/a)}
\begin{tabularx}{\textwidth}{@{}p{0.18\textwidth}X@{}}
\textbf{Priority/status} & P2 / \nolinkurl{included_acquired_research_copy} \\
\textbf{Authors} & not recorded in compact catalog \\
\textbf{DOI / identifier} & not recorded \\
\textbf{Archive path} & \path{02_MATLAB_AND_SIMULATION/additional_acquired/2026_UCSB_Introduction_to_mmWave_Radar_Sensing_Lab_Notes.pdf} \\
\end{tabularx}
\textbf{Why it is here.} Listed in prior acquisition staging; no separate PDF bytes available in this runtime. Physical research copy is included in the topical archive; check original publication license before redistribution.

\textbf{What to extract for this project.} Use to map the abstract model to AWR2944/DCA1000 configuration fields, timing windows, sample rates, clocking, calibration, channel topology, and raw-data interpretation.

\textbf{Notebook prompt.} Write down the paper's assumptions in equation form. Mark each assumption as \emph{true in V0}, \emph{true only after hardware synchronization}, \emph{requires calibration}, or \emph{not true for two independent AWR boards}. This classification is more valuable than a generic summary.

\vspace{0.35em}\hrule\vspace{0.65em}

\subsection*{59. Merlo - Picosecond NLOS Wireless Time/Frequency Transfer presentation (year n/a)}
\begin{tabularx}{\textwidth}{@{}p{0.18\textwidth}X@{}}
\textbf{Priority/status} & P2 / \nolinkurl{included_acquired_research_copy} \\
\textbf{Authors} & not recorded in compact catalog \\
\textbf{DOI / identifier} & not recorded \\
\textbf{Archive path} & \path{03_TWO_WAY_TIME_TRANSFER_AND_SYNC/additional_acquired/Merlo_Picosecond_NLOS_Wireless_Time_Frequency_Transfer_Presentation.pdf} \\
\end{tabularx}
\textbf{Why it is here.} Listed in prior acquisition staging; no separate PDF bytes available in this runtime. Physical research copy is included in the topical archive; check original publication license before redistribution.

\textbf{What to extract for this project.} Use to formalize clock epoch offset, clock-rate/frequency offset, reciprocity, two-way cancellation, presynchronization, and calibration. Translate each distinct clock/path error into an explicit simulator parameter.

\textbf{Notebook prompt.} Write down the paper's assumptions in equation form. Mark each assumption as \emph{true in V0}, \emph{true only after hardware synchronization}, \emph{requires calibration}, or \emph{not true for two independent AWR boards}. This classification is more valuable than a generic summary.

\vspace{0.35em}\hrule\vspace{0.65em}

\clearpage
\section{reference - source documentation}

\subsection*{60. 2013 Pfeffer FMCW MIMO Frequency Division TX Beamforming (year n/a)}
\begin{tabularx}{\textwidth}{@{}p{0.18\textwidth}X@{}}
\textbf{Priority/status} & reference / \nolinkurl{included_restricted_user_supplied} \\
\textbf{Authors} & not recorded in compact catalog \\
\textbf{DOI / identifier} & not recorded \\
\textbf{Archive path} & \path{09_RESTRICTED_OR_CITATION_ONLY/user_supplied_restricted/2013_Pfeffer_FMCW_MIMO_Frequency_Division_TX_Beamforming.pdf} \\
\end{tabularx}
\textbf{Why it is here.} user\_upload

\textbf{What to extract for this project.} Reference for FMCW/radar fundamentals or adjacent methods. Use it to sanity-check notation, signal relationships, and assumptions when its topic becomes active in the implementation.

\textbf{Notebook prompt.} Write down the paper's assumptions in equation form. Mark each assumption as \emph{true in V0}, \emph{true only after hardware synchronization}, \emph{requires calibration}, or \emph{not true for two independent AWR boards}. This classification is more valuable than a generic summary.

\vspace{0.35em}\hrule\vspace{0.65em}

\subsection*{61. 2022 Kumbul Smoothed Phase Coded FMCW (year n/a)}
\begin{tabularx}{\textwidth}{@{}p{0.18\textwidth}X@{}}
\textbf{Priority/status} & reference / \nolinkurl{included_topical} \\
\textbf{Authors} & not recorded in compact catalog \\
\textbf{DOI / identifier} & not recorded \\
\textbf{Archive path} & \path{04_CODED_FMCW_AND_MULTIRADAR/phase_coded_FMCW/2022_Kumbul_Smoothed_Phase_Coded_FMCW.pdf} \\
\end{tabularx}
\textbf{Why it is here.} source\_or\_downloaded

\textbf{What to extract for this project.} Use only after the uncoded V0/V1 model is validated. Extract how waveform identity is embedded, how synchronization/code alignment is performed, and whether dechirp preserves the low-bandwidth receiver advantage.

\textbf{Notebook prompt.} Write down the paper's assumptions in equation form. Mark each assumption as \emph{true in V0}, \emph{true only after hardware synchronization}, \emph{requires calibration}, or \emph{not true for two independent AWR boards}. This classification is more valuable than a generic summary.

\vspace{0.35em}\hrule\vspace{0.65em}

\subsection*{62. 2022 Merlo Mghabghab Nanzer Wireless Picosecond Time Synchronization (year n/a)}
\begin{tabularx}{\textwidth}{@{}p{0.18\textwidth}X@{}}
\textbf{Priority/status} & reference / \nolinkurl{included_topical} \\
\textbf{Authors} & not recorded in compact catalog \\
\textbf{DOI / identifier} & not recorded \\
\textbf{Archive path} & \path{03_TWO_WAY_TIME_TRANSFER_AND_SYNC/additional_acquired/2022_Merlo_Mghabghab_Nanzer_Wireless_Picosecond_Time_Synchronization.pdf} \\
\end{tabularx}
\textbf{Why it is here.} source\_or\_downloaded

\textbf{What to extract for this project.} Use to formalize clock epoch offset, clock-rate/frequency offset, reciprocity, two-way cancellation, presynchronization, and calibration. Translate each distinct clock/path error into an explicit simulator parameter.

\textbf{Notebook prompt.} Write down the paper's assumptions in equation form. Mark each assumption as \emph{true in V0}, \emph{true only after hardware synchronization}, \emph{requires calibration}, or \emph{not true for two independent AWR boards}. This classification is more valuable than a generic summary.

\vspace{0.35em}\hrule\vspace{0.65em}

\subsection*{63. 2023 Merlo High Accuracy Wireless Time Frequency Transfer Distributed Beamforming (year n/a)}
\begin{tabularx}{\textwidth}{@{}p{0.18\textwidth}X@{}}
\textbf{Priority/status} & reference / \nolinkurl{included_topical} \\
\textbf{Authors} & not recorded in compact catalog \\
\textbf{DOI / identifier} & not recorded \\
\textbf{Archive path} & \path{03_TWO_WAY_TIME_TRANSFER_AND_SYNC/picosecond_wireless_sync/2023_Merlo_High_Accuracy_Wireless_Time_Frequency_Transfer_Distributed_Beamforming.pdf} \\
\end{tabularx}
\textbf{Why it is here.} source\_or\_downloaded

\textbf{What to extract for this project.} Use for the precision/error-budget layer. Implement the impairment before the estimator whenever possible, then compare bias/variance/RMSE with the paper's analytical or measured behavior.

\textbf{Notebook prompt.} Write down the paper's assumptions in equation form. Mark each assumption as \emph{true in V0}, \emph{true only after hardware synchronization}, \emph{requires calibration}, or \emph{not true for two independent AWR boards}. This classification is more valuable than a generic summary.

\vspace{0.35em}\hrule\vspace{0.65em}

\subsection*{64. 2024 Decentralized Picosecond Synchronization Distributed Wireless Systems (year n/a)}
\begin{tabularx}{\textwidth}{@{}p{0.18\textwidth}X@{}}
\textbf{Priority/status} & reference / \nolinkurl{included_topical} \\
\textbf{Authors} & not recorded in compact catalog \\
\textbf{DOI / identifier} & not recorded \\
\textbf{Archive path} & \path{03_TWO_WAY_TIME_TRANSFER_AND_SYNC/picosecond_wireless_sync/2024_Decentralized_Picosecond_Synchronization_Distributed_Wireless_Systems.pdf} \\
\end{tabularx}
\textbf{Why it is here.} source\_or\_downloaded

\textbf{What to extract for this project.} Use to formalize clock epoch offset, clock-rate/frequency offset, reciprocity, two-way cancellation, presynchronization, and calibration. Translate each distinct clock/path error into an explicit simulator parameter.

\textbf{Notebook prompt.} Write down the paper's assumptions in equation form. Mark each assumption as \emph{true in V0}, \emph{true only after hardware synchronization}, \emph{requires calibration}, or \emph{not true for two independent AWR boards}. This classification is more valuable than a generic summary.

\vspace{0.35em}\hrule\vspace{0.65em}

\subsection*{65. 2024 Shandi Merlo Nanzer Wireless Picosecond Time Synchronization Dynamic Connectivity (year n/a)}
\begin{tabularx}{\textwidth}{@{}p{0.18\textwidth}X@{}}
\textbf{Priority/status} & reference / \nolinkurl{included_topical} \\
\textbf{Authors} & not recorded in compact catalog \\
\textbf{DOI / identifier} & not recorded \\
\textbf{Archive path} & \path{03_TWO_WAY_TIME_TRANSFER_AND_SYNC/picosecond_wireless_sync/2024_Shandi_Merlo_Nanzer_Wireless_Picosecond_Time_Synchronization_Dynamic_Connectivity.pdf} \\
\end{tabularx}
\textbf{Why it is here.} source\_or\_downloaded

\textbf{What to extract for this project.} Use to formalize clock epoch offset, clock-rate/frequency offset, reciprocity, two-way cancellation, presynchronization, and calibration. Translate each distinct clock/path error into an explicit simulator parameter.

\textbf{Notebook prompt.} Write down the paper's assumptions in equation form. Mark each assumption as \emph{true in V0}, \emph{true only after hardware synchronization}, \emph{requires calibration}, or \emph{not true for two independent AWR boards}. This classification is more valuable than a generic summary.

\vspace{0.35em}\hrule\vspace{0.65em}

\subsection*{66. 2026-08-07 Saeed Meeting Notes (year n/a)}
\begin{tabularx}{\textwidth}{@{}p{0.18\textwidth}X@{}}
\textbf{Priority/status} & reference / \nolinkurl{included_topical} \\
\textbf{Authors} & not recorded in compact catalog \\
\textbf{DOI / identifier} & not recorded \\
\textbf{Archive path} & \path{07_PROJECT_CONTEXT/2026-08-07_Saeed_Meeting_Notes.pdf} \\
\end{tabularx}
\textbf{Why it is here.} source\_or\_downloaded

\textbf{What to extract for this project.} Reference for FMCW/radar fundamentals or adjacent methods. Use it to sanity-check notation, signal relationships, and assumptions when its topic becomes active in the implementation.

\textbf{Notebook prompt.} Write down the paper's assumptions in equation form. Mark each assumption as \emph{true in V0}, \emph{true only after hardware synchronization}, \emph{requires calibration}, or \emph{not true for two independent AWR boards}. This classification is more valuable than a generic summary.

\vspace{0.35em}\hrule\vspace{0.65em}

\subsection*{67. BIPM TWSTFT Calibration Guidelines 2016 (year n/a)}
\begin{tabularx}{\textwidth}{@{}p{0.18\textwidth}X@{}}
\textbf{Priority/status} & reference / \nolinkurl{included_topical} \\
\textbf{Authors} & not recorded in compact catalog \\
\textbf{DOI / identifier} & not recorded \\
\textbf{Archive path} & \path{03_TWO_WAY_TIME_TRANSFER_AND_SYNC/foundations_and_standards/BIPM_TWSTFT_Calibration_Guidelines_2016.pdf} \\
\end{tabularx}
\textbf{Why it is here.} source\_or\_downloaded

\textbf{What to extract for this project.} Use to formalize clock epoch offset, clock-rate/frequency offset, reciprocity, two-way cancellation, presynchronization, and calibration. Translate each distinct clock/path error into an explicit simulator parameter.

\textbf{Notebook prompt.} Write down the paper's assumptions in equation form. Mark each assumption as \emph{true in V0}, \emph{true only after hardware synchronization}, \emph{requires calibration}, or \emph{not true for two independent AWR boards}. This classification is more valuable than a generic summary.

\vspace{0.35em}\hrule\vspace{0.65em}

\subsection*{68. FMCW TWTT acquisition status 2026-08-10.md (year n/a)}
\begin{tabularx}{\textwidth}{@{}p{0.18\textwidth}X@{}}
\textbf{Priority/status} & reference / \nolinkurl{included_topical} \\
\textbf{Authors} & not recorded in compact catalog \\
\textbf{DOI / identifier} & not recorded \\
\textbf{Archive path} & not present \\
\end{tabularx}
\textbf{Why it is here.} source\_or\_downloaded

\textbf{What to extract for this project.} Reference for FMCW/radar fundamentals or adjacent methods. Use it to sanity-check notation, signal relationships, and assumptions when its topic becomes active in the implementation.

\textbf{Notebook prompt.} Write down the paper's assumptions in equation form. Mark each assumption as \emph{true in V0}, \emph{true only after hardware synchronization}, \emph{requires calibration}, or \emph{not true for two independent AWR boards}. This classification is more valuable than a generic summary.

\vspace{0.35em}\hrule\vspace{0.65em}

\subsection*{69. Honeywell EP2985625A1 FMCW Radar with Timing Synchronization (year n/a)}
\begin{tabularx}{\textwidth}{@{}p{0.18\textwidth}X@{}}
\textbf{Priority/status} & reference / \nolinkurl{included_topical} \\
\textbf{Authors} & not recorded in compact catalog \\
\textbf{DOI / identifier} & not recorded \\
\textbf{Archive path} & \path{03_TWO_WAY_TIME_TRANSFER_AND_SYNC/patents/Honeywell_EP2985625A1_FMCW_Radar_with_Timing_Synchronization.pdf} \\
\end{tabularx}
\textbf{Why it is here.} source\_or\_downloaded

\textbf{What to extract for this project.} Use to formalize clock epoch offset, clock-rate/frequency offset, reciprocity, two-way cancellation, presynchronization, and calibration. Translate each distinct clock/path error into an explicit simulator parameter.

\textbf{Notebook prompt.} Write down the paper's assumptions in equation form. Mark each assumption as \emph{true in V0}, \emph{true only after hardware synchronization}, \emph{requires calibration}, or \emph{not true for two independent AWR boards}. This classification is more valuable than a generic summary.

\vspace{0.35em}\hrule\vspace{0.65em}

\subsection*{70. Honeywell US20160047892A1 FMCW Radar Phase Encoded Data Channel (year n/a)}
\begin{tabularx}{\textwidth}{@{}p{0.18\textwidth}X@{}}
\textbf{Priority/status} & reference / \nolinkurl{included_topical} \\
\textbf{Authors} & not recorded in compact catalog \\
\textbf{DOI / identifier} & not recorded \\
\textbf{Archive path} & \path{04_CODED_FMCW_AND_MULTIRADAR/phase_coded_FMCW/Honeywell_US20160047892A1_FMCW_Radar_Phase_Encoded_Data_Channel.pdf} \\
\end{tabularx}
\textbf{Why it is here.} source\_or\_downloaded

\textbf{What to extract for this project.} Reference for FMCW/radar fundamentals or adjacent methods. Use it to sanity-check notation, signal relationships, and assumptions when its topic becomes active in the implementation.

\textbf{Notebook prompt.} Write down the paper's assumptions in equation form. Mark each assumption as \emph{true in V0}, \emph{true only after hardware synchronization}, \emph{requires calibration}, or \emph{not true for two independent AWR boards}. This classification is more valuable than a generic summary.

\vspace{0.35em}\hrule\vspace{0.65em}

\subsection*{71. MathWorks Design of FMCW Radars for Active Safety Applications (year n/a)}
\begin{tabularx}{\textwidth}{@{}p{0.18\textwidth}X@{}}
\textbf{Priority/status} & reference / \nolinkurl{included_topical} \\
\textbf{Authors} & not recorded in compact catalog \\
\textbf{DOI / identifier} & not recorded \\
\textbf{Archive path} & \path{02_MATLAB_AND_SIMULATION/MathWorks/MathWorks_Design_of_FMCW_Radars_for_Active_Safety_Applications.pdf} \\
\end{tabularx}
\textbf{Why it is here.} source\_or\_downloaded

\textbf{What to extract for this project.} Reference for FMCW/radar fundamentals or adjacent methods. Use it to sanity-check notation, signal relationships, and assumptions when its topic becomes active in the implementation.

\textbf{Notebook prompt.} Write down the paper's assumptions in equation form. Mark each assumption as \emph{true in V0}, \emph{true only after hardware synchronization}, \emph{requires calibration}, or \emph{not true for two independent AWR boards}. This classification is more valuable than a generic summary.

\vspace{0.35em}\hrule\vspace{0.65em}

\subsection*{72. MathWorks Radar System Design Using MATLAB Simulink 2016 (year n/a)}
\begin{tabularx}{\textwidth}{@{}p{0.18\textwidth}X@{}}
\textbf{Priority/status} & reference / \nolinkurl{included_topical} \\
\textbf{Authors} & not recorded in compact catalog \\
\textbf{DOI / identifier} & not recorded \\
\textbf{Archive path} & \path{02_MATLAB_AND_SIMULATION/MathWorks/MathWorks_Radar_System_Design_Using_MATLAB_Simulink_2016.pdf} \\
\end{tabularx}
\textbf{Why it is here.} source\_or\_downloaded

\textbf{What to extract for this project.} Use to define simulator architecture, signal-domain fidelity, validation strategy, or efficient baseband implementation. Extract module boundaries, required state variables, and validation metrics.

\textbf{Notebook prompt.} Write down the paper's assumptions in equation form. Mark each assumption as \emph{true in V0}, \emph{true only after hardware synchronization}, \emph{requires calibration}, or \emph{not true for two independent AWR boards}. This classification is more valuable than a generic summary.

\vspace{0.35em}\hrule\vspace{0.65em}

\subsection*{73. NIST Fundamentals of Two Way Time Transfers by Satellite (year n/a)}
\begin{tabularx}{\textwidth}{@{}p{0.18\textwidth}X@{}}
\textbf{Priority/status} & reference / \nolinkurl{included_topical} \\
\textbf{Authors} & not recorded in compact catalog \\
\textbf{DOI / identifier} & not recorded \\
\textbf{Archive path} & \path{03_TWO_WAY_TIME_TRANSFER_AND_SYNC/foundations_and_standards/NIST_Fundamentals_of_Two_Way_Time_Transfers_by_Satellite.pdf} \\
\end{tabularx}
\textbf{Why it is here.} source\_or\_downloaded

\textbf{What to extract for this project.} Use to formalize clock epoch offset, clock-rate/frequency offset, reciprocity, two-way cancellation, presynchronization, and calibration. Translate each distinct clock/path error into an explicit simulator parameter.

\textbf{Notebook prompt.} Write down the paper's assumptions in equation form. Mark each assumption as \emph{true in V0}, \emph{true only after hardware synchronization}, \emph{requires calibration}, or \emph{not true for two independent AWR boards}. This classification is more valuable than a generic summary.

\vspace{0.35em}\hrule\vspace{0.65em}

\subsection*{74. TI AWR2943 AWR2944 Datasheet RevE 2026 (year n/a)}
\begin{tabularx}{\textwidth}{@{}p{0.18\textwidth}X@{}}
\textbf{Priority/status} & reference / \nolinkurl{included_topical} \\
\textbf{Authors} & not recorded in compact catalog \\
\textbf{DOI / identifier} & not recorded \\
\textbf{Archive path} & \path{06_TI_AWR2944_IMPLEMENTATION/AWR2944/TI_AWR2943_AWR2944_Datasheet_RevE_2026.pdf} \\
\end{tabularx}
\textbf{Why it is here.} source\_or\_downloaded

\textbf{What to extract for this project.} Use to map the abstract model to AWR2944/DCA1000 configuration fields, timing windows, sample rates, clocking, calibration, channel topology, and raw-data interpretation.

\textbf{Notebook prompt.} Write down the paper's assumptions in equation form. Mark each assumption as \emph{true in V0}, \emph{true only after hardware synchronization}, \emph{requires calibration}, or \emph{not true for two independent AWR boards}. This classification is more valuable than a generic summary.

\vspace{0.35em}\hrule\vspace{0.65em}

\subsection*{75. TI AWR2944 AWR2944P EVM User Guide RevC (year n/a)}
\begin{tabularx}{\textwidth}{@{}p{0.18\textwidth}X@{}}
\textbf{Priority/status} & reference / \nolinkurl{included_topical} \\
\textbf{Authors} & not recorded in compact catalog \\
\textbf{DOI / identifier} & not recorded \\
\textbf{Archive path} & \path{06_TI_AWR2944_IMPLEMENTATION/AWR2944/TI_AWR2944_AWR2944P_EVM_User_Guide_RevC.pdf} \\
\end{tabularx}
\textbf{Why it is here.} source\_or\_downloaded

\textbf{What to extract for this project.} Use to map the abstract model to AWR2944/DCA1000 configuration fields, timing windows, sample rates, clocking, calibration, channel topology, and raw-data interpretation.

\textbf{Notebook prompt.} Write down the paper's assumptions in equation form. Mark each assumption as \emph{true in V0}, \emph{true only after hardware synchronization}, \emph{requires calibration}, or \emph{not true for two independent AWR boards}. This classification is more valuable than a generic summary.

\vspace{0.35em}\hrule\vspace{0.65em}

\subsection*{76. TI AWR294x Device Errata RevC (year n/a)}
\begin{tabularx}{\textwidth}{@{}p{0.18\textwidth}X@{}}
\textbf{Priority/status} & reference / \nolinkurl{included_topical} \\
\textbf{Authors} & not recorded in compact catalog \\
\textbf{DOI / identifier} & not recorded \\
\textbf{Archive path} & \path{06_TI_AWR2944_IMPLEMENTATION/AWR2944/TI_AWR294x_Device_Errata_RevC.pdf} \\
\end{tabularx}
\textbf{Why it is here.} source\_or\_downloaded

\textbf{What to extract for this project.} Use to map the abstract model to AWR2944/DCA1000 configuration fields, timing windows, sample rates, clocking, calibration, channel topology, and raw-data interpretation.

\textbf{Notebook prompt.} Write down the paper's assumptions in equation form. Mark each assumption as \emph{true in V0}, \emph{true only after hardware synchronization}, \emph{requires calibration}, or \emph{not true for two independent AWR boards}. This classification is more valuable than a generic summary.

\vspace{0.35em}\hrule\vspace{0.65em}

\subsection*{77. TI AWR294x Technical Reference Manual RevD (year n/a)}
\begin{tabularx}{\textwidth}{@{}p{0.18\textwidth}X@{}}
\textbf{Priority/status} & reference / \nolinkurl{included_topical} \\
\textbf{Authors} & not recorded in compact catalog \\
\textbf{DOI / identifier} & not recorded \\
\textbf{Archive path} & \path{06_TI_AWR2944_IMPLEMENTATION/AWR2944/TI_AWR294x_Technical_Reference_Manual_RevD.pdf} \\
\end{tabularx}
\textbf{Why it is here.} source\_or\_downloaded

\textbf{What to extract for this project.} Use to map the abstract model to AWR2944/DCA1000 configuration fields, timing windows, sample rates, clocking, calibration, channel topology, and raw-data interpretation.

\textbf{Notebook prompt.} Write down the paper's assumptions in equation form. Mark each assumption as \emph{true in V0}, \emph{true only after hardware synchronization}, \emph{requires calibration}, or \emph{not true for two independent AWR boards}. This classification is more valuable than a generic summary.

\vspace{0.35em}\hrule\vspace{0.65em}

\subsection*{78. TI Cascade Coherency Phase Shifter Calibration (year n/a)}
\begin{tabularx}{\textwidth}{@{}p{0.18\textwidth}X@{}}
\textbf{Priority/status} & reference / \nolinkurl{included_topical} \\
\textbf{Authors} & not recorded in compact catalog \\
\textbf{DOI / identifier} & not recorded \\
\textbf{Archive path} & \path{06_TI_AWR2944_IMPLEMENTATION/chirp_config_calibration_and_MIMO/TI_Cascade_Coherency_Phase_Shifter_Calibration.pdf} \\
\end{tabularx}
\textbf{Why it is here.} source\_or\_downloaded

\textbf{What to extract for this project.} Use to map the abstract model to AWR2944/DCA1000 configuration fields, timing windows, sample rates, clocking, calibration, channel topology, and raw-data interpretation.

\textbf{Notebook prompt.} Write down the paper's assumptions in equation form. Mark each assumption as \emph{true in V0}, \emph{true only after hardware synchronization}, \emph{requires calibration}, or \emph{not true for two independent AWR boards}. This classification is more valuable than a generic summary.

\vspace{0.35em}\hrule\vspace{0.65em}

\subsection*{79. TI Fundamentals of mmWave Radar Sensors (year n/a)}
\begin{tabularx}{\textwidth}{@{}p{0.18\textwidth}X@{}}
\textbf{Priority/status} & reference / \nolinkurl{included_topical} \\
\textbf{Authors} & not recorded in compact catalog \\
\textbf{DOI / identifier} & not recorded \\
\textbf{Archive path} & \path{01_FMCW_FUNDAMENTALS/TI_and_intro/TI_Fundamentals_of_mmWave_Radar_Sensors.pdf} \\
\end{tabularx}
\textbf{Why it is here.} source\_or\_downloaded

\textbf{What to extract for this project.} Use to map the abstract model to AWR2944/DCA1000 configuration fields, timing windows, sample rates, clocking, calibration, channel topology, and raw-data interpretation.

\textbf{Notebook prompt.} Write down the paper's assumptions in equation form. Mark each assumption as \emph{true in V0}, \emph{true only after hardware synchronization}, \emph{requires calibration}, or \emph{not true for two independent AWR boards}. This classification is more valuable than a generic summary.

\vspace{0.35em}\hrule\vspace{0.65em}

\subsection*{80. TI Getting Started with mmWave Sensors 2025 (year n/a)}
\begin{tabularx}{\textwidth}{@{}p{0.18\textwidth}X@{}}
\textbf{Priority/status} & reference / \nolinkurl{included_topical} \\
\textbf{Authors} & not recorded in compact catalog \\
\textbf{DOI / identifier} & not recorded \\
\textbf{Archive path} & \path{01_FMCW_FUNDAMENTALS/TI_and_intro/TI_Getting_Started_with_mmWave_Sensors_2025.pdf} \\
\end{tabularx}
\textbf{Why it is here.} source\_or\_downloaded

\textbf{What to extract for this project.} Use to map the abstract model to AWR2944/DCA1000 configuration fields, timing windows, sample rates, clocking, calibration, channel topology, and raw-data interpretation.

\textbf{Notebook prompt.} Write down the paper's assumptions in equation form. Mark each assumption as \emph{true in V0}, \emph{true only after hardware synchronization}, \emph{requires calibration}, or \emph{not true for two independent AWR boards}. This classification is more valuable than a generic summary.

\vspace{0.35em}\hrule\vspace{0.65em}

\subsection*{81. TI MIMO Radar (year n/a)}
\begin{tabularx}{\textwidth}{@{}p{0.18\textwidth}X@{}}
\textbf{Priority/status} & reference / \nolinkurl{included_topical} \\
\textbf{Authors} & not recorded in compact catalog \\
\textbf{DOI / identifier} & not recorded \\
\textbf{Archive path} & \path{06_TI_AWR2944_IMPLEMENTATION/chirp_config_calibration_and_MIMO/TI_MIMO_Radar.pdf} \\
\end{tabularx}
\textbf{Why it is here.} source\_or\_downloaded

\textbf{What to extract for this project.} Use to map the abstract model to AWR2944/DCA1000 configuration fields, timing windows, sample rates, clocking, calibration, channel topology, and raw-data interpretation.

\textbf{Notebook prompt.} Write down the paper's assumptions in equation form. Mark each assumption as \emph{true in V0}, \emph{true only after hardware synchronization}, \emph{requires calibration}, or \emph{not true for two independent AWR boards}. This classification is more valuable than a generic summary.

\vspace{0.35em}\hrule\vspace{0.65em}

\subsection*{82. TI Programming Chirp Parameters (year n/a)}
\begin{tabularx}{\textwidth}{@{}p{0.18\textwidth}X@{}}
\textbf{Priority/status} & reference / \nolinkurl{included_topical} \\
\textbf{Authors} & not recorded in compact catalog \\
\textbf{DOI / identifier} & not recorded \\
\textbf{Archive path} & \path{06_TI_AWR2944_IMPLEMENTATION/chirp_config_calibration_and_MIMO/TI_Programming_Chirp_Parameters.pdf} \\
\end{tabularx}
\textbf{Why it is here.} source\_or\_downloaded

\textbf{What to extract for this project.} Use to map the abstract model to AWR2944/DCA1000 configuration fields, timing windows, sample rates, clocking, calibration, channel topology, and raw-data interpretation.

\textbf{Notebook prompt.} Write down the paper's assumptions in equation form. Mark each assumption as \emph{true in V0}, \emph{true only after hardware synchronization}, \emph{requires calibration}, or \emph{not true for two independent AWR boards}. This classification is more valuable than a generic summary.

\vspace{0.35em}\hrule\vspace{0.65em}

\subsection*{83. TI Self Calibration mmWave Radar Devices (year n/a)}
\begin{tabularx}{\textwidth}{@{}p{0.18\textwidth}X@{}}
\textbf{Priority/status} & reference / \nolinkurl{included_topical} \\
\textbf{Authors} & not recorded in compact catalog \\
\textbf{DOI / identifier} & not recorded \\
\textbf{Archive path} & \path{06_TI_AWR2944_IMPLEMENTATION/chirp_config_calibration_and_MIMO/TI_Self_Calibration_mmWave_Radar_Devices.pdf} \\
\end{tabularx}
\textbf{Why it is here.} source\_or\_downloaded

\textbf{What to extract for this project.} Use to map the abstract model to AWR2944/DCA1000 configuration fields, timing windows, sample rates, clocking, calibration, channel topology, and raw-data interpretation.

\textbf{Notebook prompt.} Write down the paper's assumptions in equation form. Mark each assumption as \emph{true in V0}, \emph{true only after hardware synchronization}, \emph{requires calibration}, or \emph{not true for two independent AWR boards}. This classification is more valuable than a generic summary.

\vspace{0.35em}\hrule\vspace{0.65em}

\subsection*{84. TI mmWave Sensor Raw Data Capture DCA1000 mmWave Studio (year n/a)}
\begin{tabularx}{\textwidth}{@{}p{0.18\textwidth}X@{}}
\textbf{Priority/status} & reference / \nolinkurl{included_topical} \\
\textbf{Authors} & not recorded in compact catalog \\
\textbf{DOI / identifier} & not recorded \\
\textbf{Archive path} & \path{06_TI_AWR2944_IMPLEMENTATION/DCA1000_and_capture/TI_mmWave_Sensor_Raw_Data_Capture_DCA1000_mmWave_Studio.pdf} \\
\end{tabularx}
\textbf{Why it is here.} source\_or\_downloaded

\textbf{What to extract for this project.} Use to map the abstract model to AWR2944/DCA1000 configuration fields, timing windows, sample rates, clocking, calibration, channel topology, and raw-data interpretation.

\textbf{Notebook prompt.} Write down the paper's assumptions in equation form. Mark each assumption as \emph{true in V0}, \emph{true only after hardware synchronization}, \emph{requires calibration}, or \emph{not true for two independent AWR boards}. This classification is more valuable than a generic summary.

\vspace{0.35em}\hrule\vspace{0.65em}
\clearpage
\section{Cross-cutting extraction checklist}
When the MATLAB model grows, use this checklist to decide which source class to reopen.
\begin{longtable}{p{0.27\textwidth}p{0.65\textwidth}}
\toprule Question & Source class to consult \\ \midrule
Why does the delayed chirp become a tone? & FMCW fundamentals and simulator papers. \\
How do time offset and CFO appear together? & 2007 cooperative synchronization lineage. \\
Why can two directions cancel oscillator terms? & TWTT standards and CFDDS-TWR. \\
Why is an FFT bin not the precision limit? & Scherr/CRLB/frequency-phase estimation literature. \\
How should phase noise be synthesized? & Herzel, Ayhan, Tschapek, range-correlation literature. \\
What is a realistic chirp nonlinearity model? & Ayhan, Wang, Tschapek, PLL/ramp papers. \\
How do we represent independent nodes? & Cooperative/coherent radar-network and modern synchronization papers. \\
How do we identify multiple chirps/users? & Lampel/Uysal/Kumbul/PMCW/orthogonal-waveform papers. \\
Which AWR fields set $B,T_c,S,F_s,N$? & TI chirp-programming, datasheet, TRM, EVM, DCA1000 material. \\
What does 10 ps really require? & Precision/error-bound + calibration + picosecond wireless time-transfer sources together. \\
\bottomrule
\end{longtable}

\section{Stop condition for the literature phase}
The literature phase is sufficiently mature for V0/V1 when every term in the ideal two-way model has a source and a unit test. It is sufficiently mature for a 10-ps claim only when the model contains, or explicitly bounds, clock-rate error, carrier-frequency offset, phase noise with realistic color/correlation, chirp nonlinearity, receiver/ADC noise, fixed and asymmetric group delays, estimator bias/variance, and calibration uncertainty. The latter is a later milestone; it should not delay the first Saeed plots.
\end{document}
